\documentclass[11pt]{article}

\usepackage[margin=1in]{geometry}
\usepackage{amsmath}
\usepackage{amssymb}
\usepackage{booktabs}
\usepackage{xcolor}
\usepackage{tcolorbox}
\definecolor{LnaDynamic}{HTML}{155A83}
\definecolor{LnaAlgebraic}{HTML}{974707}
\definecolor{LnaLink}{HTML}{284E6B}
\newenvironment{rowpair}[1]{%
  \def\rowcolor{#1}%
  \begin{tcolorbox}[colback=#1!3!white,colframe=#1,
    coltext=black,boxrule=0.4pt,leftrule=2pt,arc=0pt,
    left=7pt,right=7pt,top=5pt,bottom=3pt,
    before skip=7pt,after skip=7pt]
  \setlength{\abovedisplayskip}{6pt}%
  \setlength{\belowdisplayskip}{6pt}%
  \setlength{\abovedisplayshortskip}{4pt}%
  \setlength{\belowdisplayshortskip}{4pt}%
  \noindent\textbf{Model equation.}\par
}{\end{tcolorbox}\nopagebreak[4]}
\newenvironment{equationblock}{%
  \par\noindent\begin{minipage}{\linewidth}%
}{\end{minipage}\par\addvspace{12pt}}
\newenvironment{linearized}[1]{%
  \par\smallskip\begingroup\color{\rowcolor}%
  \noindent\textbf{Linearized #1.}\par
}{\endgroup}
\setcounter{tocdepth}{1}
\clubpenalty=10000
\widowpenalty=10000
\usepackage[T1]{fontenc}
\usepackage{microtype}
\usepackage[colorlinks=true,linkcolor=LnaLink,urlcolor=LnaLink]{hyperref}

\newcommand{\dd}{\mathrm{d}}
\newcommand{\e}{\mathrm{e}}
\newcommand{\code}[1]{\texttt{\detokenize{#1}}}
\newcommand{\sourcefile}[1]{\path{#1}}

\title{Radial Linear Nonadiabatic Analysis in MESA/star\\
Discretization, Convection Closures, and Diagnostics}
\author{}
\date{2026-09-19}

\begin{document}

\maketitle

\begin{abstract}
This document specifies the radial linear nonadiabatic operator in MESA/star.
Section~\ref{sec:equation-system} pairs each model
equation with its colored linearized row and source routine. It distinguishes
the velocity grids, energy forms, pressure work, and RSP2, TDC, frozen-flux,
and static MLT closures. RSP2 solves independent face superadiabaticity and
cell turbulent velocity, using shared TDC lengths. The optional three equation
model adds face entropy flux and variance to the turbulence evolution system;
the default model retains its signed algebraic convection closure.
The appendices retain closure details, work diagnostics, and schema 130 export.
\end{abstract}

\begingroup
\small
\tableofcontents
\endgroup
\clearpage

\section{Scope and conventions}

The implementation is a MESA/star operator, not a wrapper around the RSP LINA
matrix.  RSP defines the comparison output and growth rate convention.  The
matrix rows use MESA/star variables and the residual branches listed below.
MESA/star indexes the surface with $k=1$ and the innermost cell with
$k=N\equiv\texttt{s\%nz}$.  The normal mode convention is
\begin{equation}
  \delta q_k(t)=\widehat{q}_k\e^{\sigma t},
  \qquad \sigma=\gamma+\mathrm{i}\omega .
\end{equation}
A positive $\gamma$ is unstable.  For $\omega>0$, the period is
\begin{equation}
  \Pi=\frac{2\pi}{\omega}.
\end{equation}

The base perturbation variables in each zone are
\begin{equation}
  \mathbf{x}_k=
  \left(\delta\ln\rho,\,\delta\ln r,\,\delta v_{\mathrm{rad}},\,
  \delta\ln T,\,\delta L\right)_k .
\end{equation}
Here $v_{\mathrm{rad}}$ is the face variable $v$ or the cell variable $u$,
according to the active hydro momentum equation.
RSP2 adds face $\delta w$ and $\delta Y_{\rm face}$, with $e_{t,f}=w_f^2$.
The gas-cell projection is $E_{t,k}=\tfrac12(w_k^2+w_{k+1}^2)$,
so $\delta E_{t,k}=w_k\delta w_k+w_{k+1}\delta w_{k+1}$.
The mathematical inner face $N+1$ has prescribed zero turbulent energy.
Its face pressure scale height is derived, not an independent unknown.
With \code{RSP2_use_3equation_model}, RSP2 also evolves face
$\mathrm{Pi}=\langle v_r's'\rangle$ and
$\mathrm{Phi}=\langle s'^2\rangle$. These roman names denote the code variables;
they are distinct from period $\Pi$ and gravitational potential $\Phi$ below.
Active TDC adds an internal
$\delta w$ variable used only by LNA; it does not add a persistent variable to
the normal star solver. In the TDC formulas below, this internal amplitude is
written as $A=v_c/\sqrt{2/3}$ to distinguish it from RSP2 $w$.


\subsection{Equation key and notation}

The colored equations are the rows assembled into the global eigenproblem:
\begin{itemize}
\setlength{\itemsep}{2pt}
\item \textcolor{LnaDynamic}{\textbf{Blue: dynamic row.}}
      $\delta F=\sigma\,\delta Q$ contributes to both matrix sides.
\item \textcolor{LnaAlgebraic}{\textbf{Rust: algebraic row.}}
      $\delta G=0$ has a zero row on the time-derivative side.
\end{itemize}
Each box keeps the model equation and its linearization together. The row
labels also identify the type without color. Unboxed equations give
coefficients, local closures, thermodynamic comparisons, or output
quantities; they do not add independent rows to the global system.

In linearized equations, all unperturbed coefficients are evaluated at the
static background. Mass and composition are fixed. The symbol $\delta$ denotes
the complex mode amplitude; the common factor $\e^{\sigma t}$ is omitted.
For any AD quantity $f$,
\begin{equation}
  \delta f=\sum_j\left.\frac{\partial f}{\partial x_j}\right|_0\delta x_j.
  \label{eq:ad-linearization}
\end{equation}
Thus a term such as $\delta U_q$ or $\delta L_c$ includes the full supported
AD response of its closure, not only its explicit velocity dependence.
Fixed background row scales are suppressed in the displayed equations; they
multiply every term in a row and do not change its solutions.
Use $\mathcal A_k=4\pi r_k^2$ for face area and $\mathcal V_k$ for cell
volume. The convection amplitude $A$ and the matrix $\mathbf A$ are distinct.

The five base rows are density, radius, momentum, energy, and transport.
Perturbed TDC adds one convection row; RSP2 adds a turbulent-energy row and a
face-flux row, with two further moment rows in its optional three equation
model. Boundary conditions and forced-zero constraints
replace the indicated rows. They do not increase the number of unknowns.

\subsection{Independent choices in the operator}

The velocity grid and the energy form are separate choices. With
\code{u_flag}, the velocity unknown is cell-centered $u$ and the face
velocity is reconstructed. Otherwise the LNA uses a face velocity unknown
$v$, including on a background with \code{v_flag} off. This changes the
momentum stencil and the face quantities used by radius and work; it does
not select an energy form.

The \code{dedt} energy option is supported for all accepted convection paths.
RSP2 also supports \code{eps_grav}, with either velocity grid.
Within \code{dedt}, \path{use_P_d_1_div_rho_form_of_work} chooses both the
work stencil and the energy inertia:
\begin{center}
\small
\begin{tabular}{@{}lll@{}}
\toprule
Work control & Pressure work & Energy differentiated in time\\
\midrule
True & Cell $P\,\dd(1/\rho)/\dd t$ & $e+E_t$\\
False & Face $P\mathcal A v_f$ flux difference
      & $e+E_t+K+\Phi$\\
\bottomrule
\end{tabular}
\end{center}
Here $E_t$ is the selected cell turbulent energy, $K$ is specific kinetic
energy, and $\Phi$ is specific gravitational potential energy.
Section~\ref{sec:energy-row} gives the two substitutions in the same
linearized energy row. The work control acts independently of velocity
time centering. For RSP2 \code{eps_grav}, thermal storage and EOS compression
enter through the entropy coefficients; turbulent storage and non-EOS local
pressure work remain explicit. This form contains no kinetic or potential
inertia, regardless of the work-control setting. Non-RSP2 LNA still requires
\code{dedt}.

\subsection{Envelope domain}

Routine \code{setup_star_LNA_var_map} uses all $N$ cells when
\code{star_LNA_T_inner} is nonpositive. For a positive value, it retains the
contiguous surface cells $1\leq k\leq N_{\rm LNA}$ whose temperatures do not
exceed that limit. Selection stops at the first hotter cell and requires at
least two active cells. The background mesh and stellar structure are retained.

The matrix contains only active variables. Routine
\code{add_ad_partials_to_matrix} omits columns outside this domain, imposing
$\delta q_k=0$ for $k>N_{\rm LNA}$. Coefficients at the cut still use the
full-model background; the cut does not replace the physical inner boundary
at $N$ or set a new \code{R_center}. Output records both zone counts. A
velocity kick sets the excluded core velocities to zero.



\section{Equations and their linearization}
\label{sec:equation-system}

Read Sections~\ref{sec:density-row}--\ref{sec:transport-row} for the five base
rows. Then use the TDC row in Section~\ref{sec:tdc-row}, or the two RSP2 rows
in Section~\ref{sec:rsp2-rows}, when that convection model is active.
Section~\ref{sec:boundary-rows} gives the surface and forced-zero replacements.
All rows refer to the same static background and mode vector.

\begin{equationblock}
\subsection{Density: fixed cell mass}
\label{sec:density-row}

The cell volume is
$\mathcal V_k=(4\pi/3)(r_k^3-r_{k+1}^3)$.
\begin{rowpair}{LnaAlgebraic}
\begin{equation}
  \rho_k\mathcal V_k=\Delta m_k.
\end{equation}
\begin{linearized}{algebraic row}
\begin{equation}
  \delta\ln\rho_k+
  \frac{3r_k^3\delta\ln r_k-3r_{k+1}^3\delta\ln r_{k+1}}
       {r_k^3-r_{k+1}^3}=0.
  \label{eq:density}
\end{equation}
\end{linearized}
\end{rowpair}
\textit{Code:} \code{assemble_density_rows} and
\code{cell_volume_for_star_LNA} in \code{star_LNA_support.f90}.
At a fixed inner radius, its perturbation is zero. The same rule applies to
columns below a temperature-selected envelope cut.
\end{equationblock}

\begin{equationblock}
\subsection{Radius: kinematics}

The face velocity is $v_f=v$ on the face grid. With \code{u_flag}, $v_f$ is
the Riemann contact velocity reconstructed from cell velocity $u$.
\begin{rowpair}{LnaDynamic}
\begin{equation}
  \frac{\dd\ln r_k}{\dd t}=\frac{v_{f,k}}{r_k}.
\end{equation}
\begin{linearized}{dynamic row}
\begin{equation}
  \frac{\delta v_{f,k}}{r_k}=\sigma\,\delta\ln r_k.
  \label{eq:radius}
\end{equation}
\end{linearized}
\end{rowpair}
\textit{Code:} \code{assemble_radius_rows}.
The background velocity is zero, so the product of its value with
$\delta r$ vanishes. The Riemann reconstruction remains differentiated:
$\delta v_f$ generally contains pressure, density, radius, and cell-velocity
responses. It is not an arithmetic average of neighboring $\delta u$.
At the surface, the reconstruction gives $v_{f,1}=u_1$.
Appendix~\ref{app:riemann-state} gives the contact-state formulas.
For a cell constrained to zero velocity, the algebraic replacement is given
in Section~\ref{sec:forced-rows}.
\end{equationblock}

\begin{equationblock}
\subsection{Momentum: the active velocity grid}
\label{sec:momentum-row}

\paragraph{Face velocity (\code{v_flag}, or the LNA face variable).}
Write $\Delta P_k=\Delta P_{{\rm tot},k}+\Delta P_{{\rm MLT},k}$, where
$\Delta P_{{\rm tot},k}=P_{{\rm tot},k-1}-P_{{\rm tot},k}$.
The pressure closure is given in Appendix~\ref{app:pressure-energy} and
$U_q$ in Appendix~\ref{app:viscosity}.
\begin{rowpair}{LnaDynamic}
\begin{equation}
  \frac{\dd v_k}{\dd t}=-\frac{G_km_{g,k}}{r_k^2}
      +U_{q,k}-\frac{\mathcal A_k}{\Delta m_{f,k}}\Delta P_k.
  \label{eq:momentum}
\end{equation}
\begin{linearized}{dynamic row}
\begin{equation}
  \left(\frac{2G_km_{g,k}}{r_k^2}
       -\frac{2\mathcal A_k\Delta P_k}{\Delta m_{f,k}}\right)\delta\ln r_k
  +\delta U_{q,k}-\frac{\mathcal A_k}{\Delta m_{f,k}}\delta\Delta P_k
  =\sigma\,\delta v_k.
  \label{eq:momentum-linear}
\end{equation}
\end{linearized}
\end{rowpair}
\textit{Code:} \code{assemble_momentum_rows} and
\code{momentum_rhs_for_star_LNA}. The face mass is fixed:
$\Delta m_{f,k}=(\Delta m_{k-1}+\Delta m_k)/2$ in the interior and
$\Delta m_{f,1}=\Delta m_1/2$ at the surface.
\end{equationblock}

\begin{equationblock}
\paragraph{Cell velocity (\code{u_flag}).}
Define the pressure, geometry, and gravity force sum
\begin{equation}
  \mathcal F_{h,k}=P_{f,k+1}\mathcal A_{k+1}-P_{f,k}\mathcal A_k
    +P_k(\mathcal A_k-\mathcal A_{k+1})+F_{g,k}.
\end{equation}
\begin{rowpair}{LnaDynamic}
\begin{equation}
  \frac{\dd u_k}{\dd t}=\frac{\mathcal F_{h,k}}{\Delta m_k}+U_{q,k}.
  \label{eq:umomentum}
\end{equation}
\begin{linearized}{dynamic row}
\begin{equation}
  \frac{\delta\mathcal F_{h,k}}{\Delta m_k}+\delta U_{q,k}
    =\sigma\,\delta u_k.
\end{equation}
\end{linearized}
\end{rowpair}
\textit{Code:} \code{eval_Riemann_dudt_rhs} in \code{hydro_riemann.f90},
called by \code{momentum_rhs_for_star_LNA}, followed by
\code{Uq_cell_for_star_LNA}. For example,
$\delta(P_f\mathcal A)=\mathcal A\,\delta P_f+
2P_f\mathcal A\,\delta\ln r$; all force terms are differentiated this way.
Here $P_k$ is the cell pressure used by the Riemann reconstruction,
including its active turbulent-pressure terms. $P_f$ is the reconstructed
contact pressure, not an interpolation of EOS pressure alone. The gravity
source is
\begin{equation}
  F_{g,k}=-\frac{\Delta m_k}{2}
    \left(\frac{G_km_k}{r_k^2}
         +\frac{G_{k+1}m_{k+1}}{r_{k+1}^2}\right),\qquad k<N.
\end{equation}
At $k=N$, only the first gravity contribution remains, the inward pressure
flux is zero, and the geometry source is $P_N\mathcal A_N$, including when
\code{R_center > 0}. At the surface, the outward momentum flux uses
$P_{\rm bc}\mathcal A_1$. Its pressure need not equal the surface pressure
used in energy work.

The cell equation therefore changes both the force stencil and the
reconstructed quantities; replacing the symbol $v$ by $u$ in the face
equation would not give this row. Finite-step pressure and radius time
centering is excluded. The separate cell eddy-viscosity source is described
in Appendix~\ref{app:viscosity}.
\end{equationblock}

\begin{equationblock}
\subsection{Energy: luminosity divergence, sources, and work}
\label{sec:energy-row}

For \code{dedt}, both work choices use the following row, with different
$e_{\rm eff}$ and $\mathcal W$ as specified below. Let
$D_mL_k=(L_k-L_{k+1})/\Delta m_k$.
\begin{rowpair}{LnaDynamic}
\begin{equation}
  \frac{\dd e_{{\rm eff},k}}{\dd t}
    =-D_mL_k+\epsilon_k-\mathcal W_k.
  \label{eq:energy}
\end{equation}
\begin{linearized}{dynamic row}
\begin{equation}
  -D_m\delta L_k+\delta\epsilon_k-\delta\mathcal W_k
    =\sigma\,\delta e_{{\rm eff},k}.
  \label{eq:energy-linear}
\end{equation}
\end{linearized}
\end{rowpair}
\textit{Code:} \path{assemble_energy_rows},
\path{energy_rhs_for_star_LNA}, and \path{energy_inertia_for_star_LNA}.
The thermal contribution is
\begin{equation}
  \delta e_k=\rho_k\left(\frac{\partial e}{\partial\rho}\right)_T
      \delta\ln\rho_k+C_{V,k}T_k\delta\ln T_k.
  \label{eq:thermalinertia}
\end{equation}
The source $\epsilon$ contains the supported nuclear, neutrino, and fixed
heating terms. Their implementation and the energy coefficients are given
in Appendix~\ref{app:pressure-energy}. The work formulas below define
$\delta\mathcal W$ in this row; they are not additional global equations.
\end{equationblock}

\begin{equationblock}
\paragraph{Local pressure work: internal plus turbulent energy.}
For \path{use_P_d_1_div_rho_form_of_work = .true.}, define the cell volume
rate $\dot{\mathcal V}_k=\mathcal A_kv_{f,k}
-\mathcal A_{k+1}v_{f,k+1}$. Then
\begin{align}
  e_{{\rm eff},k}&=e_k+E_{t,k},&
  \mathcal W_k^{\rm local}
    &=\frac{P_{c,k}}{\Delta m_k}\dot{\mathcal V}_k,\\
  \delta e_{{\rm eff},k}&=\delta e_k+\delta E_{t,k},&
  \delta\mathcal W_k^{\rm local}
    &=\frac{P_{c,k,0}}{\Delta m_k}
       \left(\mathcal A_k\delta v_{f,k}
             -\mathcal A_{k+1}\delta v_{f,k+1}\right).
  \label{eq:simplepressurework}
\end{align}
$P_c$ is the cell work pressure from
\code{dwork_dm_for_star_LNA}: the selected EOS/RSP2 pressure plus the
MLT turbulent-pressure coefficient where enabled. The same $P_{c,k,0}$
multiplies both bounding volume fluxes. The nonlinear counterpart is
\code{hydro_energy:eval_simple_PdV_work}.

On the full domain, mass conservation and the radius rows give
\begin{equation}
  \frac{\delta\dot{\mathcal V}_k}{\Delta m_k}
   =\sigma\,\delta(1/\rho_k)
   =-\frac{\sigma}{\rho_k}\delta\ln\rho_k.
\end{equation}
Thus the work contribution can be moved to the inertia side as
$-(P_{c,k,0}/\rho_k)\delta\ln\rho_k$ after eliminating the velocity
divergence. The code retains that divergence in $\mathbf A$ and
$\delta e+\delta E_t$ in $\mathbf B$. It adds neither $\delta K$ nor
$\delta\Phi$ for this work choice.
\end{equationblock}

\begin{equationblock}
\paragraph{Face pressure work: total energy.}
For \path{use_P_d_1_div_rho_form_of_work = .false.}, the paired choices are
\begin{align}
  e_{{\rm eff},k}&=e_k+E_{t,k}+K_k+\Phi_k,\\
  \mathcal W_k^{\rm flux}
    &=\frac{P_{f,k}\mathcal A_kv_{f,k}
            -P_{f,k+1}\mathcal A_{k+1}v_{f,k+1}}{\Delta m_k},\\
  \delta e_{{\rm eff},k}
    &=\delta e_k+\delta E_{t,k}+\delta K_k+\delta\Phi_k,\\
  \delta\mathcal W_k^{\rm flux}
    &=\frac{P_{f,k,0}\mathcal A_k\delta v_{f,k}
            -P_{f,k+1,0}\mathcal A_{k+1}\delta v_{f,k+1}}{\Delta m_k}.
  \label{eq:pressurework}
\end{align}
Now the two work coefficients are different face pressures.
On the face-velocity grid, \path{static_face_pressure_for_star_LNA} forms
them from mass-weighted cell pressures and the selected MLT contribution.
With \code{u_flag}, they are the background Riemann pressures
\path{s%P_face_ad}. At the surface of the cell-velocity scheme,
$P_{f,1}=P_{{\rm eos},1}$ for energy work, even when the momentum
boundary uses a different $P_{\rm bc}$.

\code{mechanical_energy_inertia_for_star_LNA} adds the kinetic and
potential derivatives only for \code{dedt} with this work control false.
On a static background $\delta K=0$, but $\delta\Phi$ is first order
in the radius perturbations and must remain. Appendix~\ref{app:mechanical}
gives the discrete energies for both velocity grids.

Changing the work form therefore changes both $\mathcal W$ and
$e_{\rm eff}$. It is not a row rescaling, and the two discrete operators
need not give identical modes. The pressure-gradient work represented by
the face flux is paired with the mechanical energy derivatives.
\end{equationblock}

\begin{equationblock}
\paragraph{Static limit, time centering, and boundaries.}
Both work forms use the same face velocity: the independent face variable
for the face scheme, and the differentiated Riemann contact velocity for
\code{u_flag}. Holding the work pressure at its background value does not
freeze that contact-velocity response. At $v_{f,0}=0$,
\begin{equation}
  \delta(P\mathcal A v_f)=P_0\mathcal A_0\,\delta v_f.
\end{equation}
Terms $\delta P\,\mathcal A_0v_{f,0}$ and
$P_0\delta\mathcal A\,v_{f,0}$ vanish. Eddy-viscosity heating and the
mechanical power $vU_q$ are also second order, although
$\delta U_q$ remains in momentum.

The LNA uses continuous time derivatives. It does not differentiate a
finite timestep or retain the nonlinear velocity/pressure time weights.
The choice of local versus flux work still applies; turning on nonlinear
velocity time centering does not select the local work form.

At the physical inner boundary $k=N$, the fixed inward velocity has
zero perturbation, so the inward term in $\delta\mathcal W_N$ is absent.
A temperature-selected envelope cut instead retains the full-model
stencil and discards excluded variable columns. It does not introduce
a new physical boundary or replace its pressure coefficients.
\end{equationblock}

\begin{equationblock}
\paragraph{RSP2 \code{eps_grav}: entropy and turbulent storage.}
This is an alternative energy row for RSP2 with either velocity grid.
Define the separate non-EOS work by
\begin{equation}
 \mathcal W_{\neg\mathrm{EOS},k}
 =\frac{P_{c,k}-P_{{\rm eos},k}}{\Delta m_k}\dot{\mathcal V}_k.
\end{equation}
The EOS compression is already contained in $\epsilon_{\rm grav}$.
\begin{rowpair}{LnaDynamic}
\begin{equation}
 \dot E_{t,k}=-D_mL_k+\epsilon_k+\epsilon_{{\rm grav},k}
                    -\mathcal W_{\neg\mathrm{EOS},k}.
\end{equation}
\begin{linearized}{dynamic row}
\begin{align}
 -D_m\delta L_k+\delta\epsilon_k-\delta\mathcal W_{\neg\mathrm{EOS},k}
 &=\sigma\bigg\{\delta E_{t,k}\notag\\
 &\quad+T_kC_{P,k}\Big[(1-\nabla_{{\rm ad},k}\chi_{T,k})\delta\ln T_k
       -\nabla_{{\rm ad},k}\chi_{\rho,k}\delta\ln\rho_k\Big]\bigg\}.
 \label{eq:epsgrav-row}
\end{align}
\end{linearized}
\end{rowpair}
This display uses the standard thermal coefficients with
\code{eps_grav_factor = 1}; Appendix~\ref{app:epsgrav} gives the implemented
latent-heat and entropy-blend substitutions. The non-EOS work linearization
uses the background $P_c-P_{\rm eos}$ times the perturbed volume rate.
\textit{Code:} \code{energy_inertia_for_star_LNA} and
\code{dwork_dm_for_star_LNA}. The work control does not select a face-work
or mechanical-inertia branch for this energy option.
\end{equationblock}

\begin{equationblock}
\paragraph{Relation to local-work \code{dedt}.}
For a consistent EOS at fixed composition, write $\nu=1/\rho$. Then
\begin{align}
 \epsilon_{\rm grav}&=-\dot e-P_{\rm eos}\dot\nu,\\
 \delta\epsilon_{\rm grav}
 &=-\sigma\left(\delta e-\frac{P_{\rm eos}}{\rho}\delta\ln\rho\right).
 \label{eq:epsgrav-comparison}
\end{align}
Using mass conservation in the local-work \code{dedt} row gives this same
continuous thermodynamic relation, plus the turbulent terms.
$\epsilon_{\rm grav}$ is not simply $-\dot\Phi$. Adding explicit
EOS pressure work or $\dot e$ to the \code{eps_grav} row would count those
terms twice. Numerical EOS derivatives and finite-step time discretizations
can differ, so this relation does not guarantee identical numerical modes or
nonlinear timesteps. The two formulations remain separately documented.
\end{equationblock}

\begin{equationblock}
\subsection{Transport: one row in the luminosity slot}
\label{sec:transport-row}

Routine \code{assemble_luminosity_rows} selects exactly one of the following
interior rows. The surface replacement is in
Section~\ref{sec:thermal-boundary}. The matrix unknown in this slot is
$\delta L$, even when the equation is written as a temperature gradient.
The nonlinear solver retains the existing name \code{equL} for this transport
row. RSP2's additional flux row is \code{rsp2_flux}, with index
\code{i_rsp2_flux}; row labels do not assign one equation to one unknown.

\paragraph{Perturbed TDC or frozen convective flux.}
The face closure supplies $L_r$ and $L_c$.
\begin{rowpair}{LnaAlgebraic}
\begin{equation}
  L_{r,k}+L_{c,k}-L_k=0.
  \label{eq:tdclum}
\end{equation}
\begin{linearized}{algebraic row}
\begin{equation}
  \delta L_{r,k}+\delta L_{c,k}-\delta L_k=0.
  \label{eq:tdclum-linear}
\end{equation}
\end{linearized}
\end{rowpair}
For \code{frozen_flux}, replace $L_c$ by its fixed background value
$L_{c,0}$ and set $\delta L_c=0$ in this same row.
\textit{Code:} \code{tdc_luminosity_resid_for_star_LNA} or
\code{frozen_flux_luminosity_resid_for_star_LNA}.
The TDC coefficients are given in Appendix~\ref{app:tdc}.
\end{equationblock}

\begin{equationblock}
\paragraph{Static MLT or RSP2: logarithmic-pressure form.}
Here $P=P_{\rm eos}$. Define
$\Delta\ln P_k=\ln P_{k-1}-\ln P_k$ and similarly for $T$.
This branch has priority when
\path{use_gradT_actual_vs_gradT_MLT_for_T_gradient_eqn} is true.
\begin{rowpair}{LnaAlgebraic}
\begin{equation}
  \nabla_k\Delta\ln P_k-\Delta\ln T_k=0.
  \label{eq:tgrad}
\end{equation}
\begin{linearized}{algebraic row}
\begin{equation}
  \Delta\ln P_k\,\delta\nabla_k
    +\nabla_k\,\delta\Delta\ln P_k-\delta\Delta\ln T_k=0.
\end{equation}
\end{linearized}
\end{rowpair}
\textit{Code:} \code{temperature_gradient_resid_for_star_LNA}.
The EOS supplies $\delta\ln P=\chi_\rho\delta\ln\rho+
\chi_T\delta\ln T$. For RSP2,
$\nabla=\nabla_{L,\rm eff}+Y_{\rm face}$ and
$\delta\nabla=\delta\nabla_{L,\rm eff}+\delta Y_{\rm face}$.
The independent flux row in Section~\ref{sec:rsp2-rows} supplies
$L_r+L_c+L_t=L$; no local inversion for $Y_{\rm face}$ remains.
\end{equationblock}

\begin{equationblock}
\paragraph{Static MLT or RSP2: radiation-pressure form.}
If the preceding branch is off and
\path{use_dPrad_dm_form_of_T_gradient_eqn} is true, define
\begin{equation}
  K_f=\frac{\Delta m_f\kappa_f}{c\mathcal A_f^2\lambda_f},
  \qquad \Delta P_{{\rm rad},k}=\frac{a}{3}(T_{k-1}^4-T_k^4).
\end{equation}
\begin{rowpair}{LnaAlgebraic}
\begin{equation}
  -K_fL_{r,k}-\Delta P_{{\rm rad},k}=0.
  \label{eq:dpraddm}
\end{equation}
\begin{linearized}{algebraic row}
\begin{equation}
  -K_f\delta L_{r,k}-L_{r,k}\delta K_f
  -\frac{4a}{3}\left(T_{k-1}^4\delta\ln T_{k-1}
                    -T_k^4\delta\ln T_k\right)=0.
\end{equation}
\end{linearized}
\end{rowpair}
\textit{Code:} \code{dPrad_dm_resid_for_star_LNA}.
At fixed mass, $\delta K_f/K_f=\delta\kappa_f/\kappa_f-
4\delta\ln r_f-\delta\lambda_f/\lambda_f$.
Appendix~\ref{app:transport} specifies the opacity floor, radiative flux limiter,
radiative luminosity, and inner-boundary spacing.
\end{equationblock}

\begin{equationblock}
\paragraph{Static MLT or RSP2: QHSE form.}
When both preceding switches are off, define
\begin{equation}
  h_f=(\dd\ln P/\dd m)_{\rm QHSE},\qquad
  \Delta m_{c,f}=\frac{\Delta m_{k-1}+\Delta m_k}{2},
\end{equation}
with the mass-weighted temperature
\begin{equation}
  T_{*,f}=a_fT_k+(1-a_f)T_{k-1},\qquad
  a_f=\frac{\Delta m_{k-1}}{\Delta m_{k-1}+\Delta m_k}.
\end{equation}
\begin{rowpair}{LnaAlgebraic}
\begin{equation}
  \Delta m_{c,f}h_f\nabla_k-\frac{T_{k-1}-T_k}{T_{*,f}}=0.
\end{equation}
\begin{linearized}{algebraic row}
\begin{equation}
\begin{aligned}
  0={}&\Delta m_{c,f}(h_f\delta\nabla_k+\nabla_k\delta h_f)
       -\frac{\delta T_{k-1}-\delta T_k}{T_{*,f}}\\
     &+\frac{T_{k-1}-T_k}{T_{*,f}^2}\,\delta T_{*,f}.
\end{aligned}
\end{equation}
\end{linearized}
\end{rowpair}
\textit{Code:} \code{temperature_gradient_resid_for_star_LNA} and
\code{star_LNA_eval_dlnPdm_qhse}. Here $\delta T=T\delta\ln T$ and
$\delta T_{*,f}=a_f\delta T_k+(1-a_f)\delta T_{k-1}$.
For active TDC, the selected spatial relation instead supplies the gradient
inside its closure; the global row remains equation~\eqref{eq:tdclum-linear}.
\end{equationblock}

\begin{equationblock}
\subsection{Additional row for perturbed TDC}
\label{sec:tdc-row}

The internal convection amplitude is $A=v_c/\sqrt{2/3}$ at the face.
Let $F_A=S_0-A D_{R0}-A^2D_0$, with the buoyancy and damping coefficients
specified in Appendix~\ref{app:tdc}. The compressive term uses face density.
\begin{rowpair}{LnaDynamic}
\begin{equation}
  2\dot A-\frac{2}{3}\alpha_{Pt}A\,\dot{\ln\rho_f}
    =S_0-A D_{R0}-A^2D_0.
\end{equation}
\begin{linearized}{dynamic row}
\begin{equation}
\begin{aligned}
  \delta S_0-(D_{R0}+2AD_0)\delta A-A\delta D_{R0}-A^2\delta D_0
  =\sigma\left(2\delta A-\frac{2}{3}\alpha_{Pt}A\delta\ln\rho_f\right).
\end{aligned}
  \label{eq:tdcvelocity}
\end{equation}
\end{linearized}
\end{rowpair}
\textit{Code:} \code{assemble_tdc_w_row},
\code{tdc_relation_for_star_LNA}, and \code{set_TDC_LNA}.
The AD inertia is represented locally by
\begin{equation}
  I_A=2A+\frac{2}{3}\alpha_{Pt}A_0\frac{\rho_{f,0}}{\rho_f},
  \qquad \delta I_A=2\delta A-\frac{2}{3}\alpha_{Pt}A_0\delta\ln\rho_f.
  \label{eq:tdcinertia}
\end{equation}
Only the variation of this expression is used; $A_0$ and $\rho_{f,0}$ are
held fixed. Nonconvective or forced-zero faces use the algebraic constraint
in Section~\ref{sec:forced-rows}. Static MLT and frozen-flux treatments do
not add this convection unknown.
\end{equationblock}

\begin{equationblock}
\subsection{Additional rows for RSP2}
\label{sec:rsp2-rows}

\paragraph{Turbulent energy.}
Use face $e_t=w^2$, $C=S-D-D_r$, and center transport luminosity $L_{t,c}$.
Let $\mu_k=(\Delta m_{k-1}+\Delta m_k)/2$ be the face mass and define
\begin{equation}
 W_{f,k}=\frac{\alpha_p}{3\mu_k}\sum_{j=k-1,k}\Delta m_j
 [\theta_P\rho_jw_k^2+(1-\theta_P)\rho_{j,0}w_{k,0}^2]
 (\rho_j^{-1}-\rho_{j,0}^{-1}).
 \label{eq:rsp2facework}
\end{equation}
Here $0$ denotes accepted state. This is a work increment, not a new
variable. The continuous work rate is its zero-step limit.
\begin{rowpair}{LnaDynamic}
\begin{equation}
 \frac{\mathrm d(w_k^2)}{\mathrm d t}+\dot W_{f,k}
 =C_k+Eq_{f,k}-\frac{L_{t,c,k-1}-L_{t,c,k}}{\mu_k}.
\end{equation}
\begin{linearized}{dynamic row}
\begin{equation}
 \begin{aligned}
 \delta C_k-\frac{\delta L_{t,c,k-1}-\delta L_{t,c,k}}{\mu_k}
 =\sigma\left[2w_k\delta w_k-
 \frac{\alpha_pw_k^2}{3\mu_k}\sum_{j=k-1,k}\Delta m_j\delta\ln\rho_j\right].
 \end{aligned}
 \label{eq:rsp2energy}
\end{equation}
\end{linearized}
\end{rowpair}
\textit{Code:} \code{assemble_rsp2_w_row},
\code{rsp2_turbulent_energy_rhs_for_star_LNA},\\
\code{rsp2_pressure_inertia_for_star_LNA}.
The pressure derivative enters the inertia matrix; this is the continuous
linearization of the nonlinear density-increment work. It does not change
implicitness or a time weight. Static viscous heating has $\delta Eq=0$.
Forced faces use $\delta w=0$. The RSP3 closure has $C=S-D$ because its
radiative cooling acts in the entropy moments.

The gas energy rows in Section~\ref{sec:energy-row} use $E_{t,k}$ and the
additional work redistribution
\begin{equation}
 \mathcal W_k\ \longleftarrow\ \mathcal W_k+
 \frac{\tfrac12(W_{f,k}+W_{f,k+1})-P_{t,c,k}^{\theta}\Delta(1/\rho_k)}{\Delta t}.
 \label{eq:rsp2redistribution}
\end{equation}
Its mass integral is zero. Both work forms and \code{eps_grav} use it.
In the linear rows its added inertia is the half sum of the face density
terms in equation~\eqref{eq:rsp2energy}, plus
$(P_{t,c,k}/\rho_k)\delta\ln\rho_k$ to subtract the native cell contribution.
Thus the displayed native gas work formulas retain their differences,
with this same conservative redistribution added for RSP2.
\end{equationblock}

\begin{equationblock}
\paragraph{Independent face superadiabaticity and flux balance.}
The new face unknown is signed $Y_{{\rm face},k}$, stored in \code{i_Y}.
The gradient and its variation are
\begin{equation}
 \nabla_k=\nabla_{L,\rm eff,k}+Y_{{\rm face},k},\qquad
 \delta\nabla_k=\delta\nabla_{L,\rm eff,k}+\delta Y_{{\rm face},k}.
\end{equation}
\begin{rowpair}{LnaAlgebraic}
\begin{equation}
 \begin{cases}
 L_{r,k}+L_{c,k}+L_{t,k}-L_k=0,&k>1,\\
 Y_{{\rm face},1}=0,&k=1.
 \end{cases}
\end{equation}
\begin{linearized}{algebraic row}
\begin{equation}
 \begin{cases}
 \delta L_{r,k}+\delta L_{c,k}+\delta L_{t,k}-\delta L_k=0,&k>1,\\
 \delta Y_{{\rm face},1}=0,&k=1.
 \end{cases}
 \label{eq:rsp2flux}
\end{equation}
\end{linearized}
\end{rowpair}
\textit{Code:} \code{assemble_rsp2_flux_row} and
\code{hydro_rsp2:rsp2_flux_residual}. The nonlinear interior row is
\begin{equation}
 R_{\mathrm{flux},k}=\frac{L_{r,k}+L_{c,k}+L_{t,k}-L_k}
 {\max\!\left(|L_{\mathrm{start},k}|,
 10^{-3}\max_j|L_{\mathrm{start},j}|,\ 1\,\mathrm{erg\,s^{-1}}\right)}.
\end{equation}
This follows the TDC luminosity scale, adding a finite floor and using absolute
values in the global maximum. Before solver iterations, current L initializes
the scale. During iterations it uses L\_start, so its derivatives are zero.
LNA uses the same residual and holds its background normalization fixed.
The surface constraint $Y_1=0$ is dimensionless. The Newton correction measure
uses $|\Delta Y|/\max(1,|Y_{\rm start}|)$: the native dimensionless column
scale with unit correction weight. This gives an absolute scale near zero
and a relative scale at large magnitude. Y is included in the norm and remains signed.
The physical surface
luminosity/temperature condition remains in its separate row. Radiative cells
retain the flux equation and may have negative Y.
\end{equationblock}

\subsection{Optional RSP2 entropy moment rows}
\label{sec:rsp2-moment-rows}

These two face rows are present only in the three equation model.
The turbulent velocity is on the same face, with $e_{t,f}=w_f^2$. The face entropy flux is
$\mathrm{Pi}_f=\langle v_r's'\rangle$, and $\mathrm{Phi}_f=\langle s'^2\rangle$
is the full entropy variance. The entropy differential and face coefficients
are specified in Appendix~\ref{app:rsp2-moments}.

\begin{equationblock}
\paragraph{Entropy flux.} Pi carries its own time dependence.
\begin{rowpair}{LnaDynamic}
\begin{equation}
\begin{aligned}
 \dot{\mathrm{Pi}}_f={}&\frac23e_{t,f}(-\partial_r s)_f
 +\left(-\frac{\partial_r P}{\rho}\right)_f
       \frac{\chi_{T,f}}{\chi_{\rho,f}C_{P,f}}\mathrm{Phi}_f\\
 &-\left(\alpha_{\mathrm{Pi}}\frac{\sqrt{e_{t,f}}}{\Lambda_f}
       +\tau_{{\rm rad},f}^{-1}+ (\partial_r v_{\rm rad})_f\right)\mathrm{Pi}_f .
\end{aligned}
\end{equation}
\begin{linearized}{dynamic row}
\begin{equation}
\begin{aligned}
 \sigma\delta\mathrm{Pi}_f={}&\frac23\left[
     (-\partial_r s)_f\delta e_{t,f}+e_{t,f}\delta(-\partial_r s)_f\right]\\
 &+\left(-\frac{\partial_r P}{\rho}\right)_f
       \frac{\chi_{T,f}}{\chi_{\rho,f}C_{P,f}}\delta\mathrm{Phi}_f\\
 &+\mathrm{Phi}_f\delta\left[
       -\frac{\partial_r P}{\rho}\frac{\chi_T}{\chi_\rho C_P}\right]_f\\
 &-\left(\alpha_{\mathrm{Pi}}\frac{\sqrt{e_{t,f}}}{\Lambda_f}
       +\tau_{{\rm rad},f}^{-1}\right)\delta\mathrm{Pi}_f\\
 &-\mathrm{Pi}_f\delta\left[
       \alpha_{\mathrm{Pi}}\frac{\sqrt{e_{t,f}}}{\Lambda_f}
       +\tau_{{\rm rad},f}^{-1}+ (\partial_r v_{\rm rad})_f\right].
\end{aligned}
\end{equation}
\end{linearized}
\end{rowpair}
The linearized expression assumes a static background. Mean strain vanishes
there, but its velocity perturbation is retained.
\end{equationblock}

\begin{equationblock}
\paragraph{Entropy variance.} Phi is the full variance of specific entropy.
\begin{rowpair}{LnaDynamic}
\begin{equation}
\begin{aligned}
 \dot{\mathrm{Phi}}_f={}&2(-\partial_r s)_f\mathrm{Pi}_f\\
 &-\left(\alpha_{\mathrm{Phi}}\frac{\sqrt{e_{t,f}}}{\Lambda_f}
             +2\tau_{{\rm rad},f}^{-1}\right)\mathrm{Phi}_f .
\end{aligned}
\end{equation}
\begin{linearized}{dynamic row}
\begin{equation}
\begin{aligned}
 \sigma\delta\mathrm{Phi}_f={}&2(-\partial_r s)_f\delta\mathrm{Pi}_f
     +2\mathrm{Pi}_f\delta(-\partial_r s)_f\\
 &-\left(\alpha_{\mathrm{Phi}}\frac{\sqrt{e_{t,f}}}{\Lambda_f}
             +2\tau_{{\rm rad},f}^{-1}\right)\delta\mathrm{Phi}_f\\
 &-\mathrm{Phi}_f\delta\left[
       \alpha_{\mathrm{Phi}}\frac{\sqrt{e_{t,f}}}{\Lambda_f}
             +2\tau_{{\rm rad},f}^{-1}\right].
\end{aligned}
\end{equation}
\end{linearized}
\end{rowpair}
\textit{Code:} \code{hydro_rsp2:rsp2_moment_rhs},
\code{do1_rsp2_moment_eqns}, and \code{assemble_rsp2_moment_rows}.
Both moments use backward Euler in nonlinear hydro and identity inertia in
LNA. The continuous LNA coefficients omit finite timestep centering.
There is no spatial transport or additional viscosity in these two rows.
The existing $L_t$ and $Eq/Uq$ terms remain in the energy and momentum equations.
\end{equationblock}

Forced faces use $\mathrm{Pi}=\mathrm{Phi}=0$. Fully dormant faces with both
local turbulent energy and both moments zero use the same zero branch.
A face with finite local turbulent energy retains both evolution rows.
LNA fixes $\delta w=0$ at a fully dormant face. At zero represented face
energy, including underflow of $w^2$, a local row without a source seed
or kinetic transport also uses this constraint when $S_k/w_k\leq0$.
The factored source ratio determines the branch on the current static
background; accepted timestep history does not enter LNA.
Positive driving or nonlocal coupling retains the physical energy row.
The w and moment boundary masks use the same face location.
Active moments at exactly zero face energy have no regular linearization
in these variables and cause an explicit LNA error.
The dormant branch does not describe spontaneous turbulence onset.

\begin{equationblock}
\paragraph{Derived scale heights and mixing lengths.}
\label{sec:cell-mixing-length}
RSP2 uses TDC's public face helpers, \code{get_TDC_Hp_face} and
\code{get_TDC_mixing_length_face}. Both models obtain cell lengths from
\code{get_TDC_mixing_length_cell}. Face lengths retain the ordinary or
reconstructed face state. The cell length uses its own EOS pressure and density,
with gravity averaged over its bounding faces:
\begin{align}
 g_{\mathrm{cell},k}&=\tfrac12(g_k+g_{k+1}),
 &g_j&=\frac{G_jm_{\mathrm{grav},j}}{r_j^2},\\
 H_{P,\mathrm{cell},k}^{\mathrm{HSE}}
   &=\frac{P_{\mathrm{eos},k}}{\rho_k g_{\mathrm{cell},k}},
 &r_{\mathrm{cell},k}&=\left[\tfrac12(r_k^3+r_{k+1}^3)\right]^{1/3}.
\end{align}
At a full centre, $g_{N+1}=0$; an excised boundary uses
$g_{N+1}=G_NM_{\rm centre}/R_{\rm centre}^2$.
For $\beta_h=\code{harmonic_dissipation_length_beta}\leq0$, the alternate-height
option replaces $H_P^{\mathrm{HSE}}$ by
$\min(H_P^{\mathrm{HSE}},\sqrt{P_{\mathrm{eos}}/G_k}/\rho)$ in that cell.
For $\beta_h>0$, the HSE height is used directly. Writing the selected height
as $H_{P,\mathrm{cell}}$, the shared length law gives
\begin{equation}
 \Lambda_{\mathrm{cell},k}=
 \begin{cases}
  \alpha_{\mathrm{MLT}}H_{P,\mathrm{cell},k},&\beta_h\leq0,\\
  \left[(\alpha_{\mathrm{MLT}}H_{P,\mathrm{cell},k})^{-1}
        +(\beta_h r_{\mathrm{cell},k})^{-1}\right]^{-1},&\beta_h>0.
 \end{cases}
 \label{eq:cell-mixing-length}
\end{equation}
The EOS state is not interpolated through faces a second time. These lengths
are differentiated functions of the cell state and bounding radii; there is
no Hp row or independent $\delta H_P$. RSP2 face w and TDC face convective
velocity retain their separate storage and viscosity averaging.
\end{equationblock}

\begin{equationblock}
\subsection{Surface rows and fixed perturbations}
\label{sec:boundary-rows}
\label{sec:thermal-boundary}
\label{sec:forced-rows}

\paragraph{Surface momentum or pressure.}
\path{assemble_surface_velocity_row} replaces the interior momentum row at
$k=1$. A fixed surface velocity has first priority and imposes
$\delta v_{{\rm rad},1}=0$, as in Section~\ref{sec:forced-rows}.
Otherwise, any of \path{use_momentum_outer_BC},
\path{use_zero_Pgas_outer_BC}, or \path{use_fixed_Psurf_outer_BC}
selects the dynamic surface equation.
\begin{rowpair}{LnaDynamic}
\begin{equation}
  \dot v_{{\rm rad},1}=F_{{\rm surf},1}(P_{\rm bc},\ldots).
\end{equation}
\begin{linearized}{dynamic row}
\begin{equation}
  \delta F_{{\rm surf},1}=\sigma\,\delta v_{{\rm rad},1}.
\end{equation}
\end{linearized}
\end{rowpair}
\textit{Code:} \code{surface_momentum_rhs_for_star_LNA}.
The force is the selected momentum formula above, with outer pressure
$P_{\rm bc}$; on the face grid, $\Delta P_{{\rm tot},1}=P_{\rm bc}-P_{{\rm tot},1}$.
\end{equationblock}

\begin{equationblock}
Without a surface momentum condition, the replacement is algebraic:
\begin{rowpair}{LnaAlgebraic}
\begin{equation}
  \ln P_{\rm bc}-\ln P_{{\rm eos},1}=0.
\end{equation}
\begin{linearized}{algebraic row}
\begin{equation}
  \delta\ln P_{\rm bc}-\chi_{\rho,1}\delta\ln\rho_1
    -\chi_{T,1}\delta\ln T_1=0.
\end{equation}
\end{linearized}
\end{rowpair}
\textit{Code:} \code{surface_P_bc_for_star_LNA} supplies either a fixed
pressure or the atmosphere pressure and its supported AD response.
\end{equationblock}

\begin{equationblock}
\paragraph{Surface luminosity or temperature.}
The surface transport slot also contains one row. For RSP2,
\code{RSP2_use_L_eqn_at_surface} selects its own luminosity boundary first.
Otherwise \code{use_RSP_L_eqn_outer_BC} selects the shared boundary, also used
with TDC. Both give the following continuous luminosity condition:
\begin{rowpair}{LnaAlgebraic}
\begin{equation}
  L_1=L_{\rm surf},\qquad
  L_{\rm surf}=f_{L,\rm surf}\,4\pi r_1^2caT_1^4.
\end{equation}
\begin{linearized}{algebraic row}
\begin{equation}
  \delta L_1-L_{\rm surf}(2\delta\ln r_1+4\delta\ln T_1)=0.
\end{equation}
\end{linearized}
\end{rowpair}
\textit{Code:} \code{rsp_lsurf_resid_for_star_LNA}, with
$f_{L,\rm surf}=\code{RSP2_Lsurf_factor}$.
\par
Otherwise the atmosphere temperature supplies the row:
\begin{rowpair}{LnaAlgebraic}
\begin{equation}
  \ln T_{\rm bc}-\ln T_1=0.
\end{equation}
\begin{linearized}{algebraic row}
\begin{equation}
  \delta\ln T_{\rm bc}-\delta\ln T_1=0.
\end{equation}
\end{linearized}
\end{rowpair}
\textit{Code:} \code{surface_temperature_resid_for_star_LNA} and
\code{surface_lnT_bc_for_star_LNA}. The atmospheric face-to-center offset
and its variation are given in Appendix~\ref{app:surface-details}.
\end{equationblock}

\begin{equationblock}
\paragraph{Forced-zero rows.}
These replace the radius row in a zero-velocity cell, the momentum row for a
fixed surface velocity, or the TDC/RSP2 convection row where the amplitude
is forced to zero. Here $z$ is that constrained variable.
\begin{rowpair}{LnaAlgebraic}
\begin{equation}
  z=z_{\rm fixed}.
\end{equation}
\begin{linearized}{algebraic row}
\begin{equation}
  \delta z=0.
\end{equation}
\end{linearized}
\end{rowpair}
\textit{Code:} \code{assemble_radius_rows},
\code{assemble_surface_velocity_row}, \code{assemble_tdc_w_row}, and
\code{assemble_rsp2_w_row}. At the envelope cut, excluded perturbations are
zero by omission of their columns, not by adding further constraint rows.
\end{equationblock}

\clearpage
\section{From the linearized rows to the eigensolve}

Routine \code{setup_star_LNA_problem} constructs a \code{star_LNA_problem}
containing the variable map, equation identifier for every row, and dense
matrices.  The top level sequence in \code{do_star_LNA} is
\begin{enumerate}
\item refresh the MESA variables with \code{set_vars_if_needed};
\item reject unsupported model branches in \code{check_star_LNA_model};
\item rebuild the Riemann face state for \code{u_flag}; RSP2 viscosity is
      a cell momentum source, not a contact-velocity shift;
\item construct the problem and report the selected variables;
\item assemble density, radius, momentum, energy, luminosity, RSP2, and TDC
      rows in that order;
\item write matrix and closure audits when requested;
\item solve and write mode products; and
\item release the problem through \code{free_star_LNA_problem}.
\end{enumerate}

The colored rows above assemble the single problem
\begin{equation}
  \mathbf{A}\widehat{\mathbf{x}}
  =\sigma\mathbf{B}\widehat{\mathbf{x}} .
  \label{eq:generalized}
\end{equation}
For $\dot q=F(q)$, the perturbation $\delta F$ enters $\mathbf{A}$ and
$\delta q$ enters $\mathbf{B}$.  Algebraic equations enter only
$\mathbf{A}$. Automatic differentiation components are mapped to neighboring
zone variables by \path{add_ad_partials_to_matrix}.

The equation registry distinguishes density, radius, forced zero velocity,
interior and surface momentum variants, energy, each luminosity closure, RSP2
turbulent energy and face flux, and active or forced zero TDC velocity rows.  It is
mechanical metadata: row formulas remain in their direct row assemblers.


\subsection{Algebraic elimination and eigenpair refinement}

Routine \code{scale_star_LNA_matrix_for_solver} scales each row and then each
column by its largest absolute entry across $\mathbf{A}$ and $\mathbf{B}$.  A
nominally dynamic variable remains in the generalized problem only when its
row contains a time derivative.  Thus a forced $\delta w_k=0$ row is
algebraic, while a TDC turbulent-velocity row remains dynamic.  Let $x_a$
denote algebraic variables and $x_d$ dynamic variables.  Since algebraic rows
have no $\mathbf{B}$ entries,
\begin{equation}
  A_{aa}x_a+A_{ad}x_d=0,
  \qquad x_a=-A_{aa}^{-1}A_{ad}x_d.
\end{equation}
Substitution gives
\begin{equation}
  \left(A_{dd}-A_{da}A_{aa}^{-1}A_{ad}\right)x_d
  =\sigma
  \left(B_{dd}-B_{da}A_{aa}^{-1}A_{ad}\right)x_d.
  \label{eq:schur}
\end{equation}
\code{build_reduced_star_LNA_problem} uses equilibrated \code{DGESVX} with
iterative refinement to evaluate $A_{aa}^{-1}A_{ad}$.  Elimination changes the
row and column norms.  Routine
\code{scale_star_LNA_pencil} therefore applies alternating row and column
max-norm scaling to the reduced matrices.  With $t_{\rm dyn}$ equal to the
model dynamical timescale, the scaled problem is
\begin{equation}
  \widetilde A y=\lambda\widetilde B y,
  \qquad
  \widetilde A=R A_{\rm red}C,
  \qquad
  \widetilde B=\frac{R B_{\rm red}C}{t_{\rm dyn}},
  \qquad
  \lambda=\sigma t_{\rm dyn}.
\end{equation}
This keeps radial-mode eigenvalues near unity during \code{DGGEV}.  The code
restores $x_d=Cy$, reconstructs $x_a$, and then reverses the full-pencil column
scaling.

Each candidate is checked against the original full matrices using the
componentwise backward error
\begin{equation}
  \epsilon_{\rm eig}=
  \max_i\frac{|(Ax-\sigma Bx)_i|}
  {\sum_j|A_{ij}x_j|+|\sigma|\sum_j|B_{ij}x_j|}.
\end{equation}
Let $\epsilon_{\rm max}=\code{star_LNA_max_eigenvector_residual}$, which
must be positive and defaults to $10^{-3}$. The triangle inequality bounds
$\epsilon_{\rm eig}$ by one; evaluation enforces this upper bound after
checking for nonfinite values. A limit of one therefore provides no
eigenvector accuracy requirement. A candidate above
$\min(\epsilon_{\rm max},10^{-1})$ enters \code{refine_star_LNA_eigenpair}.
The Newton update can change frequency when either
$\epsilon_{\rm eig}\leq10^{-1}$ or its balanced normwise residual satisfies
\begin{equation}
 \epsilon_{\rm norm}=
 \frac{\max_i|(A_s x_s-\sigma B_s x_s)_i|}
 {\max_i\left(\sum_j|(A_s)_{ij}(x_s)_j|
       +|\sigma|\sum_j|(B_s)_{ij}(x_s)_j|\right)}\leq10^{-6}.
\end{equation}
Here $A_s,B_s,x_s$ use the full matrix scaling before algebraic elimination.
Tiny moment components can have order-unity relative errors even when the
eigenpair has a small normwise defect. This second criterion allows those
components to be refined; it does not replace the final componentwise test.
On the scaled full matrices, define
$M=A-\sigma B$ and $r=Mx$. The Newton correction satisfies
\begin{equation}
  M\,\delta x-Bx\,\delta\sigma=-r,\qquad\delta x_j=0,
\end{equation}
where $j$ is the largest component of the current eigenvector at each
iteration. Two solves give
\begin{equation}
  y=M^{-1}(-r),\qquad z=M^{-1}Bx,\qquad
  \delta\sigma=-y_j/z_j,\qquad\delta x=y+z\,\delta\sigma.
\end{equation}
If neither initial criterion holds, keep $\sigma$ fixed and replace $x$
by $z/z_j$. This inverse iteration repairs the eigenvector before allowing
a frequency correction. Its trial update is a weighted average of the old
and repaired vectors, preserving tiny components without subtractive loss.
Each iteration uses \code{ZGBTRF} and \code{ZGBTRS} on the full banded
matrix. At most four refinement steps are accepted; backtracking requires a
decrease in the scaled residual. The refined pair is retained only if it also
reduces the residual against the original unscaled matrices. The residual
limit and mode filters are then applied.

The solve computes the full
reduced eigensystem.  The requested mode count limits output, not eigensolve
cost.

The scaling passes require quadratic work in the matrix dimension. Candidate
residuals use the matrix bandwidth, as does the local Newton correction.
Dense algebraic elimination and \code{DGGEV} require cubic
work and determine the large-model scaling.  A targeted sparse or iterative
solver is therefore required to make the requested mode count reduce the main
solve cost.


\clearpage
\appendix
\numberwithin{equation}{section}
\section{Pressure, energy, and face closures}
\label{app:pressure-energy}

These quantities supply coefficients and derivatives to the rows in
Section~\ref{sec:equation-system}. They are not additional global equations.

\subsection{Pressure}

The cell pressure closure shared by the operator and pressure work output is
\begin{equation}
  P_{\mathrm{tot}}=P_{\mathrm{eos}}+P_{\mathrm{t,RSP2}},
  \qquad
  P_{\mathrm{t,RSP2}}
  =\alpha_p\frac{2}{3}\rho e_t,
  \qquad e_t=E_{t,k}=\tfrac12(w_k^2+w_{k+1}^2),
  \label{eq:rsp2pressure}
\end{equation}
when RSP2 turbulent pressure perturbations are enabled.  TDC/MLT turbulent
pressure remains face centered and enters separately through
\code{d_mlt_Pturb_face_for_star_LNA}:
\begin{equation}
  P_{\mathrm{t,MLT},f}
  =f_{\mathrm{Pt}}\frac{\rho_f v_c^2}{3}.
  \label{eq:mltpressure}
\end{equation}
The separate MLT pressure term in face momentum uses one convection
velocity coefficient multiplying a cell-density difference:
\begin{align}
  \Delta P_{{\rm MLT},k}
    &=\frac{f_{\rm Pt}}{3}v_{c,k}^2(\rho_{k-1}-\rho_k),\\
  \delta\Delta P_{{\rm MLT},k}
    &=\frac{f_{\rm Pt}}{3}\left[
       v_{c,k}^2(\rho_{k-1}\delta\ln\rho_{k-1}
                       -\rho_k\delta\ln\rho_k)\right.\nonumber\\
    &\hspace{38mm}\left.
       +2v_{c,k}(\rho_{k-1}-\rho_k)\delta v_{c,k}\right].
\end{align}
These expressions apply when turbulent-pressure perturbations are enabled.
Active TDC supplies $\delta v_c$ from its internal amplitude. Otherwise the
helper uses a positive \code{mlt_vc_old} as a fixed coefficient, so
$\delta v_c=0$ in this term. The term is zero when its controls or velocity
condition are inactive. It is not a difference of independently evaluated
turbulent pressures at two adjacent faces.

\subsection{Work coefficients on the two grids}
\label{app:work-coefficients}

Section~\ref{sec:energy-row} gives both work stencils and their
linearizations. Their pressure coefficients are constructed differently.
For local work, \code{dwork_dm_for_star_LNA} uses
\begin{equation}
  P_{c,k,0}=P_{{\rm tot},k,0}
     +f_{\rm Pt}\frac{\rho_{f,k,0}v_{c,{\rm old},k}^2}{3},
\end{equation}
where the second term is included only for $k>1$,
\code{mlt_Pturb_factor > 0}, and \code{mlt_vc_old(k) > 0}.
This face-density contribution to the cell work coefficient is evaluated
by \code{get_rho_face}; it is not replaced by the Riemann pressure.

For flux work on the face-velocity grid,
\code{static_face_pressure_for_star_LNA} gives
\begin{equation}
  P_{f,k,0}=a_kP_{{\rm tot},k,0}
       +(1-a_k)P_{{\rm tot},k-1,0}+P_{{\rm t,MLT},f,k,0},
  \quad
  a_k=\frac{\Delta m_{k-1}}{\Delta m_{k-1}+\Delta m_k}.
\end{equation}
At $k=1$, the cell-pressure part is $P_{{\rm tot},1,0}$.
The MLT addition uses current $v_c$ and the TDC face state for an active
TDC face with positive velocity; otherwise it uses the positive
\code{mlt_vc_old} and \code{get_rho_face_val} when enabled.
For flux work with \code{u_flag}, the code instead takes the values of
\path{s%P_face_ad}, including the surface convention stated in the main
text. These distinctions must be retained when comparing energy and
momentum pressure terms.

\subsection{Mechanical energy inertia}
\label{app:mechanical}

For the supported case without mass corrections,
\code{star_utils:cell_specific_KE_qp} uses
\begin{equation}
  K_k=\begin{cases}
    u_k^2/2,&\code{u_flag},\\
    (v_k^2+v_{k+1}^2)/4,&\code{v_flag},\\
    0,&\text{neither hydro velocity flag}.
  \end{cases}
\end{equation}
Consequently,
\begin{equation}
  \delta K_k=\begin{cases}
    u_{k,0}\delta u_k,&\code{u_flag},\\
    (v_{k,0}\delta v_k+v_{k+1,0}\delta v_{k+1})/2,
       &\code{v_flag},\\
    0,&\text{otherwise}.
  \end{cases}
\end{equation}
Both hydrodynamic expressions vanish on a static background. The face
scheme averages the squares of the bounding velocities; it does not
square their average. At $N$, the inward velocity is fixed.

For $k<N$, \code{cell_specific_PE_qp} gives the same gravitational
energy for both velocity grids:
\begin{align}
  \Phi_k&=-\frac{1}{2}
      \left(\frac{G_km_{g,k}}{r_k}
           +\frac{G_{k+1}m_{g,k+1}}{r_{k+1}}\right),\\
  \delta\Phi_k&=\frac{1}{2}
      \left(\frac{G_km_{g,k}}{r_k}\delta\ln r_k
           +\frac{G_{k+1}m_{g,k+1}}{r_{k+1}}
             \delta\ln r_{k+1}\right).
\end{align}
This is the average of the two face potentials, not a cell-center
$-Gm/r$ approximation. At $N$, a nonzero fixed
\code{R_center} contributes a constant inward potential; its perturbation
is zero. The contribution is omitted when \code{R_center = 0}.
An envelope cut drops excluded perturbation columns and keeps the
background coefficients.

\code{mechanical_energy_inertia_for_star_LNA} calls
\code{get_dke_dt_dpe_dt} with a unit time interval to recover these
derivatives. It places them in $\mathbf B$ only when
\code{energy_eqn_option = 'dedt'} and the local-work control is false.
A retained kinetic derivative at a small nonzero background velocity
does not make the rest of the operator a nonstatic linearization.

\subsection{Turbulent energy and the RSP2 gas-energy relation}
\label{app:turbulent-energy}

The control \path{star_LNA_perturb_turbulent_energy} gates the cell
turbulent inertia. When it is false,
\path{turbulent_energy_inertia_for_star_LNA} returns zero.
Otherwise RSP2 contributes
\begin{equation}
  E_{t,k}=w_k^2,\qquad \delta E_{t,k}=2w_k\delta w_k.
\end{equation}
For TDC, inclusion additionally requires
\code{star_LNA_include_tdc}, \code{MLT_option = 'TDC'}, and
\path{TDC_include_eturb_in_energy_equation}. The cell inertia is
\begin{align}
  E_{t,k}&=\frac{3}{4}(v_{c,k}^2+v_{c,k+1}^2)
          =\frac{1}{2}(A_k^2+A_{k+1}^2),\\
  \delta E_{t,k}&=A_k\delta A_k+A_{k+1}\delta A_{k+1},
  \qquad k<N.
\end{align}
Only the first face term remains at the physical inner cell $N$.
Active TDC uses its internal LNA amplitude through
\code{turbulent_tdc_energy_inertia_for_star_LNA}. The nonactive TDC
branch, when these inclusion controls still apply, takes derivatives
from \path{s%mlt_vc_ad}; it does not introduce a new amplitude unknown.
These are cell energy coefficients, distinct from the TDC face-amplitude
inertia $I_A$ in equation~\eqref{eq:tdcinertia}.

For the local-work RSP2 case with turbulent energy and pressure included,
and without an additional MLT pressure term, the two equations are
\begin{align}
  \dot e+\dot E_t
    &=-D_m(L_r+L_c+L_t)+\epsilon+Eq
      -(P_{\rm eos}+P_t)\dot{\mathcal V}/\Delta m,\\
  \dot E_t
    &=C+Eq-D_mL_t-P_t\dot{\mathcal V}/\Delta m.
\end{align}
Subtracting the second from the first gives
\begin{align}
  \dot e&=-D_m(L_r+L_c)+\epsilon-C
            -P_{\rm eos}\dot{\mathcal V}/\Delta m,\\
  \sigma\delta e&=-D_m(\delta L_r+\delta L_c)+\delta\epsilon-\delta C
            -\frac{P_{{\rm eos},0}}{\Delta m}
               \delta\dot{\mathcal V}.
\end{align}
Thus $L_t$, $Eq$, and turbulent-pressure work cancel in the derived gas row;
none is counted twice.
This gas-energy equation is a derived combination of assembled rows,
not an extra row. In the flux-work branch the mechanical terms and
face-pressure stencil must also be retained in the subtraction.
Freezing the turbulent inertia or pressure removes the corresponding
cancellation.

\subsection{The \code{eps_grav} coefficients and their limits}
\label{app:epsgrav}

For a thermodynamically consistent EOS at fixed composition, the
thermal part of the first-law relation can also be written
\begin{equation}
  \epsilon_{\rm grav}
    =-TC_P\left[
       (1-\nabla_{\rm ad}\chi_T)\dot{\ln T}
       -\nabla_{\rm ad}\chi_\rho\dot{\ln\rho}\right].
\end{equation}
The standard MESA implementation is
\code{eps_grav:do_std_eps_grav}. Its static thermal variation is
\begin{equation}
  \delta\epsilon_{\rm grav}
    =-\sigma TC_P\left[
       (1-\nabla_{\rm ad}\chi_T)\delta\ln T
       -\nabla_{\rm ad}\chi_\rho\delta\ln\rho\right].
\end{equation}
The identity linking this to the internal-energy derivatives is
\begin{equation}
  \delta e-\frac{P_{\rm eos}}{\rho}\delta\ln\rho
    =TC_P\left[
       (1-\nabla_{\rm ad}\chi_T)\delta\ln T
       -\nabla_{\rm ad}\chi_\rho\delta\ln\rho\right].
\end{equation}
This is a thermodynamic comparison. Numerical EOS quantities, especially
across blends, need not satisfy the identity exactly. Finite-step
\code{dedt} and \code{eps_grav} residuals also have different time
discretizations. Equality of these continuous expressions therefore does
not establish identical numerical operators or finite-step solutions.

For RSP2, \code{energy_inertia_for_star_LNA} implements the two thermal
coefficients directly. Its standard branch adds the stored latent-heat terms:
\begin{align}
 B_{\ln\rho}&=-TC_P\nabla_{\rm ad}\chi_\rho
                  +\code{latent_ddlnRho},\\
 B_{\ln T}&=TC_P(1-\nabla_{\rm ad}\chi_T)
                  +\code{latent_ddlnT}.
\end{align}
In the PC EOS entropy transition it uses the same background Gamma thresholds
as \path{eval1_eps_grav_and_partials}, blending this standard expression with
$T_0\delta s$. The thermal contribution is multiplied by
\path{eps_grav_factor}; $2w\delta w$ is then added without that factor.
The blend weight is a background coefficient: its perturbation multiplies a
zero background time derivative. Composition perturbations vanish in the
fixed-composition LNA. The nonlinear time-centering weights are not part of
the continuous LNA operator. User \code{other_eps_grav} hooks are rejected.

The corresponding nonlinear RSP2 path calls the existing
\code{eval_eps_grav_and_partials}, retaining its composition and finite-step
choices. See \sourcefile{star/private/eps_grav.f90} and
\sourcefile{star/defaults/controls.defaults}. Non-RSP2 LNA retains its
\code{dedt}-only restriction.

\subsection{Source derivatives}

\code{energy_sources_for_star_LNA} retains the stored nuclear derivatives
with respect to $\ln\rho$ and $\ln T$. Nuclear heating is fixed when the
selected split-burn branch uses \code{burn_avg_epsnuc}, and the explicit
nuclear term is zero for zero \code{eps_nuc_factor} or the nonlocal
Ni/Co decay branch. The nonnuclear neutrino coefficient retains the
nonlinear half weight:
\begin{equation}
  \delta\epsilon_{\nu,{\rm nonnuc}}
   =\frac{1}{2}\left[
      \frac{\partial\epsilon_{\nu,{\rm nonnuc}}}{\partial\ln\rho}
          \delta\ln\rho+
      \frac{\partial\epsilon_{\nu,{\rm nonnuc}}}{\partial\ln T}
          \delta\ln T\right].
\end{equation}
It enters $\delta\epsilon$ with a minus sign.
\code{extra_heat} and \code{irradiation_heat} are fixed coefficients.
The diffusion, sedimentation, premixing, and phase-separation
\code{others} terms of the nonlinear energy equation supply no
additional LNA derivatives. Static viscous heating and $vU_q$ have zero
first-order variation, as described in Appendix~\ref{app:viscosity}.

\subsection{Riemann face state for \code{u_flag}}
\label{app:riemann-state}

Routine \code{hydro_riemann:do1_uface_and_Pface} uses the inner cell
$L=k$ and outer cell $R=k-1$. Let $\widetilde P_L,\widetilde P_R$
be their cell pressures after the gravity reconstruction, with
\begin{align}
  h_g&=-\frac{G_km_{g,k}}{r_k^2\mathcal A_k},\\
  \widetilde P_L&=P_L+\frac{\Delta m_k}{2}h_g,\qquad
  \widetilde P_R=P_R-\frac{\Delta m_{k-1}}{2}h_g.
\end{align}
The acoustic speeds use the unreconstructed cell pressures:
$c_{s,L}^2=\Gamma_{1,L}P_L/\rho_L$ and similarly for $R$. The selected
wave speeds and contact state are
\begin{align}
  S_L&=\min(u_L-c_{s,L},u_R-c_{s,R}),\\
  S_R&=\max(u_R+c_{s,R},u_L+c_{s,L}),\\
  Z&=\rho_R(S_R-u_R)+\rho_L(u_L-S_L),\\
  v_f=S_*&=\frac{\rho_Ru_R(S_R-u_R)+\rho_Lu_L(u_L-S_L)
                      +\widetilde P_L-\widetilde P_R}{Z},\\
  P_f&=\frac{1}{2}\left[
      \widetilde P_L+\rho_L(u_L-S_L)(u_L-S_*)
      +\widetilde P_R+\rho_R(u_R-S_R)(u_R-S_*)\right].
\end{align}
Writing the numerator of $S_*$ as $N_*$,
\begin{equation}
  \delta v_f=\frac{\delta N_*-v_{f,0}\delta Z}{Z_0}.
\end{equation}
AD differentiates the selected wave-speed branches, the EOS response,
and the gravity reconstruction. At $v_{f,0}=0$, this reduces to
$\delta N_*/Z_0$, which still contains the pressure-difference response.
The same contact velocity enters the radius equation and both work forms.
$P_f$ is differentiated in momentum; in static work only its background
value multiplies $\delta v_f$.

At $k=1$, the routine sets $v_f=u_1$ and $P_f=P_{{\rm eos},1}$
directly. The outward momentum flux separately uses $P_{\rm bc}$.
\code{do_star_LNA} rebuilds this face state. RSP2 now puts its viscous
force in cell momentum, using \code{compute_Uq_dm_cell}; it does not add
an acceleration to the contact velocity. The static counterpart is
\code{Uq_cell_for_star_LNA}. Active turbulent pressure with \code{u_flag}
still requires \path{star_LNA_perturb_turbulent_pressure}, so the momentum,
work, and reconstructed face states use consistent pressure responses.

\subsection{Radiative transport}
\label{app:transport}

The face temperature $T_{*,f}$ and spacing $\Delta m_{c,f}$ are defined in
Section~\ref{sec:transport-row}. The selected static MLT/RSP2 transport row
and its linearization are given there in full.

The spatial gradient used by active TDC is
\begin{equation}
  \nabla_{\mathrm{actual},f}=
  \frac{(T_{k-1}-T_k)/T_{*,f}}
       {\Delta m_{c,f}(\dd\ln P/\dd m)_{\rm QHSE}}.
  \label{eq:tdcqhse}
\end{equation}
Routine \code{star_LNA_eval_dlnPdm_qhse} supplies the QHSE derivative,
including the selected turbulent-pressure terms. Without those terms it is
$-Gm_f/(4\pi r_f^4P_f)$. The helper starts from the reconstructed EOS pressure
when face reconstruction is active, then adds the selected MLT turbulent
pressure. The TDC radiative conductance $L_{0,f}$ uses the reconstructed EOS
pressure. Without the turbulent-pressure addition, the same EOS pressure
therefore enters both expressions.

For the radiation-pressure branch, equation~\eqref{eq:dpraddm} gives the
assembled row. The following coefficients define that relation.

The opacity uses \code{min_kap_for_dPrad_dm_eqn} as a floor. The face state is
reconstructed when \code{use_face_reconstruction} is active. In an MLT cell,
$L_{\mathrm{rad}}$ is obtained from the background temperature gradient rather
than from the total luminosity. With flux limiting,
\begin{equation}
  R_f=\frac{\mathcal A_f|T_{k-1}^4-T_k^4|}
  {\Delta m_f\kappa_f(T_{k-1}^4+T_k^4)/2},
  \qquad
  \lambda_f=\frac{6+3R_f}{6+(3+R_f)R_f}.
\end{equation}
Without flux limiting, $\lambda_f=1$.

In \code{dPrad_dm_resid_for_star_LNA}, $\Delta m_f$ is normally
\path{s%dm_bar(k)}. At $k=N$ with \code{R_center > 0}, it is replaced by
$\Delta m_{c,f}$, matching the cell-center spacing beside an excised core.

For active TDC, this transport relation supplies the spatial gradient while the
single algebraic row remains equation~\eqref{eq:tdclum}. The implementation
first obtains
\begin{equation}
  L_{\mathrm{rad},f}=
  -\frac{c\mathcal A_f^2\lambda_f}{\Delta m_f\kappa_f}
  \frac{a}{3}\left(T_{k-1}^4-T_k^4\right),
  \qquad
  \nabla_{\mathrm{actual},f}=\frac{L_{\mathrm{rad},f}}{L_{0,f}},
  \label{eq:tdcdprad}
\end{equation}
and passes this gradient to \code{set_TDC_LNA}. Thus the TDC luminosity and
turbulent-velocity relations use the selected spatial transport equation
without adding another unknown or residual row. RSP2 instead inserts its
reconstructed $L_r$ directly in equation~\eqref{eq:dpraddm}.
The TDC helper \code{actual_gradT_for_star_LNA} still uses
\path{s%dm_bar(k)} in this branch; the explicit inner-spacing replacement
above applies to the static MLT/RSP2 row. These paths require separate
inner-boundary checks.


\section{TDC coefficients and reconstructed faces}
\label{app:tdc}

The global luminosity and amplitude rows are
in Sections~\ref{sec:transport-row} and \ref{sec:tdc-row}.
This appendix defines their local coefficients and AD inputs.

The active TDC relation is evaluated by \code{set_TDC_LNA}. The MESA/star
routine \path{tdc_face_state_for_star_LNA} obtains
\begin{equation}
  (T,\rho,P,e,C_P,\chi_\rho,\chi_T,\nabla_{\mathrm{ad}},
  \kappa,H_P,\nabla_{\mathrm{rad}})_f
\end{equation}
from \code{get_reconstructed_face_state_ad}.  When
\code{use_face_reconstruction} is true, these are the cached reconstructed EOS
and opacity quantities. Otherwise they are the
stored face values.  They remain automatic differentiation objects in both
paths.

The routine \code{actual_gradT_for_star_LNA} supplies
$\nabla_{\mathrm{actual}}$. It uses the logarithmic-pressure or QHSE spatial
relation normally and equation~\eqref{eq:tdcdprad} when the model selects
\path{use_dPrad_dm_form_of_T_gradient_eqn}. The QHSE path uses reconstructed
face pressure when face reconstruction is active.

The internal velocity variable is
\begin{equation}
  A=\frac{v_c}{\sqrt{2/3}}.
\end{equation}
Define
\begin{align}
  L_0 &=\frac{16\pi ac}{3}\frac{GmT^4}{P\kappa},
  &\nabla_r&=\frac{L}{L_0},\\
  Y&=\nabla_{\mathrm{actual}}-\nabla_{L,\rm eff}.
\end{align}
The default is $\nabla_{L,\rm eff}=\nabla_L$. With
\code{TDC_use_dynamical_gradL}, \code{get_TDC_dynamical_gradL} supplies
the following instantaneous spatial relation, using its actual local
\code{pressure_gradient_factor}:
\begin{align}
 \mathrm{pressure\_gradient\_factor}_k
   &=\frac{P_{k-1}-P_k}{\Delta P_{{\rm QHSE},k}},\\
 \nabla_{L,\rm eff,k}
   &=\mathrm{pressure\_gradient\_factor}_k\,\nabla_{L,k}.
  \label{eq:dynamicgradl}
\end{align}
Here $\Delta P_{{\rm QHSE},k}=-G_km_{g,k}\Delta m_f/(4\pi r_k^4)$
on the nonrotating, fixed-mass background accepted by LNA, with
$\Delta m_f=(\Delta m_{k-1}+\Delta m_k)/2$.
Both TDC and the default one equation RSP2 gradient closure retain the product rule:
\begin{align}
 \delta\nabla_{L,\rm eff,k}
   &=\mathrm{pressure\_gradient\_factor}_k\,\delta\nabla_{L,k}\notag\\
   &\quad+\nabla_{L,k}\,\delta\mathrm{pressure\_gradient\_factor}_k,\\
 \delta\mathrm{pressure\_gradient\_factor}_k
   &=\frac{\delta P_{k-1}-\delta P_k}{\Delta P_{{\rm QHSE},k}}
     +4\,\mathrm{pressure\_gradient\_factor}_k\,\delta\ln r_k.
\end{align}
The composition-gradient contribution in $\nabla_L$ is a fixed background
input to these AD derivatives. In RSP2, \code{compute_RSP2_gradT} stores this
response in \code{gradT_ad}; the LNA temperature and flux rows use it directly.
The optional RSP2 moment model instead uses the row consistent mapping in
Appendix~\ref{app:rsp2-moments}.
The pressure-ratio perturbation is retained even when its background value
is one. The control defaults to false. The helper leaves $\nabla_L$ unchanged
at the surface, without an active hydro velocity variable, or when the ratio
is undefined. It does the same at $k=N$ with \code{R_center > 0} if the ratio
is nonpositive.
Its pressure difference includes the selected MLT turbulent pressure, using
\code{mlt_vc_old} in that coefficient. Its Ledoux support is incomplete: the uncommon
\code{get_brunt_B} fallback already uses the QHSE pressure interval, so its
composition contribution receives an extra pressure-ratio factor.
The TDC length passed to \code{set_TDC_LNA} is
\begin{align}
  H_P^* &=
  \begin{cases}
  H_P,&\beta_h\leq 0,\\
  P/(\rho g),&\beta_h>0,
  \end{cases}\\
  \Lambda_0 &= \alpha_{\mathrm{MLT}}H_P^*,\\
  \Lambda &=
  \begin{cases}
  \Lambda_0,&\beta_h\leq 0,\\
  \left(\Lambda_0^{-1}+(\beta_h r)^{-1}\right)^{-1},&\beta_h>0,
  \end{cases}\\
  \alpha_{\mathrm{eff}} &=
  \begin{cases}
  \alpha_{\mathrm{MLT}},&\beta_h\leq 0,\\
  \Lambda/H_P^*,&\beta_h>0.
  \end{cases}
\end{align}
Here \code{get_mlt_mixing_length} evaluates $\Lambda$. Thus the LNA and
nonlinear TDC paths differentiate the same harmonic length.
The local MLT call supplies $\Gamma$.  The environmental excess is
\begin{equation}
  Y_{\mathrm{env}}=
  \begin{cases}
  Y\Gamma/(1+\Gamma),&Y>0\text{ and the MLT correction is enabled},\\
  Y,&\text{otherwise}.
  \end{cases}
  \label{eq:yenv}
\end{equation}
The TDC call to \code{set_TDC_LNA} sets the nonlinear enthalpy flux limiter
to false. The TDC coefficients are
\begin{align}
  S_0 &=\alpha_Sx_S\alpha_{\mathrm{eff}}\frac{C_PT}{H_P^*}
       \nabla_{\mathrm{ad}}Y_{\mathrm{env}}+E_q/w,
       \label{eq:tdcsource}\\
  D_0 &=\frac{\alpha_Dx_D}{\Lambda},\\
  D_{R0} &=4\sigma_{\mathrm{SB}}
    \left(\frac{\alpha_Rx_R}{\Lambda}\right)^2
    \frac{T^3}{\rho^2C_P\kappa}.
\end{align}
Here $E_q/w=0$ for the static first order operator.  The two algebraic
luminosities are
\begin{align}
  L_{\mathrm{rad}} &=L_0\nabla_{\mathrm{actual}},\\
  L_{\mathrm{conv}} &=\alpha_C\alpha_{\mathrm{eff}}\alpha_c\rho TC_P
     4\pi r^2A Y_{\mathrm{env}},
  \qquad \alpha_c=\frac{1}{2}\sqrt{\frac{2}{3}}.
\end{align}
They enter Equation~\eqref{eq:tdclum}.

The TDC cell eddy-viscosity length is the shared cell-local
$\Lambda_{\mathrm{cell},k}$ in Equation~\eqref{eq:cell-mixing-length}.
Both nonlinear TDC and \code{tdc_chi_coefficient_for_star_LNA} call
\code{get_TDC_mixing_length_cell}. At a static zero-shear background,
the perturbation of the stress coefficient multiplies zero shear, so the
first-order viscous force uses its background value times the perturbed
velocity shear. Face lengths used by the TDC convection closure and
u-grid viscosity retain their face definition.


\section{RSP2 convection coefficients}
\label{app:rsp2}

These coefficients supply the transport, turbulent-energy, and independent
flux rows in Sections~\ref{sec:transport-row} and \ref{sec:rsp2-rows}.
Their AD values are calculated in \code{hydro_rsp2.f90} and reused by
\path{star_LNA_turbulence_closures.f90}.
The PII, convective luminosity and buoyant source formulas below describe
the default one equation model. Appendix~\ref{app:rsp2-moments} gives their
replacements for the optional moment model. Mixing lengths, turbulent
pressure, kinetic energy transport and viscosity retain their existing forms.

\subsection{Face state and velocity--entropy correlation}

RSP2 stores w on faces, colocated with Y and, in RSP3, Pi and Phi.
Its face energy is $w_f^2$ and cell energy is
$E_{t,k}=\tfrac12(w_k^2+w_{k+1}^2)$.
Thermodynamic interpolation retains
\begin{equation}
 a_f=\begin{cases}
 \Delta m_{k-1}/(\Delta m_{k-1}+\Delta m_k),&\text{mass interpolation},\\
 1/2,&\text{otherwise}.
 \end{cases}
\end{equation}
These weights do not average the independent face turbulent variables.

\begin{equationblock}
PII has units of specific entropy, with the turbulent-speed factor applied
separately in luminosity and source. With $x_S=x_C=\tfrac12\sqrt{2/3}$,
\path{compute_PII_from_Y} evaluates
\begin{equation}
 \mathrm{PII}_f=x_S\frac{\Lambda_f}{Hp_f}C_{P,f}Y_f,
 \qquad Y_f=\nabla_f-\nabla_{L,\rm eff,f}.
 \label{eq:pii}
\end{equation}
Hp and $\Lambda$ come from the shared TDC length functions. The ratio
$\Lambda/Hp$ cannot in general be replaced by $\alpha_{\rm MLT}$ when
harmonic lengths are selected. PII is zero at prescribed nonturbulent faces. The inner mathematical
boundary is $N+1$; stored face N is active unless selected by the inner mask.

The same signed PII enters convective luminosity and buoyant driving.
Its perturbation is the product rule
\begin{align}
 \delta\mathrm{PII}_f
 &=x_S\frac{\Lambda_f}{Hp_f}C_{P,f}\,\delta Y_f\notag\\
 &\quad+x_SY_f\left[
   \frac{C_{P,f}}{Hp_f}\delta\Lambda_f
   +\frac{\Lambda_f}{Hp_f}\delta C_{P,f}
   -\frac{\Lambda_f C_{P,f}}{Hp_f^2}\delta Hp_f\right].
\end{align}
This form remains regular at $Y_f=0$ and retains the response for either
sign of Y. Nonlinear hydro and RSP2 LNA use the same \code{PII_ad}; face
weights are held fixed during the perturbation. There is no local inversion
of luminosity to recover Y.
\end{equationblock}

\subsection{Radiative, convective, and turbulent fluxes}

The gradient is now evaluated directly:
\begin{equation}
 \nabla_f=\nabla_{L,\rm eff,f}+Y_f,\qquad
 \delta\nabla_f=\delta\nabla_{L,\rm eff,f}+\delta Y_f.
\end{equation}
\code{compute_Lrad_coeff}, \code{compute_Lc_terms}, and \code{compute_Lt}
then supply
\begin{align}
 L_{r,f}&=\mathrm{Lrad\_coeff}_f\nabla_f,
 &\mathrm{Lrad\_coeff}_f&=
 \mathcal A_f\frac{4acT_f^4}{3\kappa_f\rho_f Hp_f},\\
 L_{c,f}&=\mathcal A_fw_f\frac{x_C}{x_S}(\rho T)_f\mathrm{PII}_f,
 \label{eq:rsp2lc}\\
 L_{t,c,k}&=-\alpha_t\mathcal A_{c,k}^2\rho_k^2\Lambda_{c,k}
 \frac{w_k+w_{k+1}}2\frac{w_k^2-w_{k+1}^2}{\Delta m_k},\\
 L_{t,f,k}&=\frac{\Delta m_kL_{t,c,k-1}+\Delta m_{k-1}L_{t,c,k}}
 {\Delta m_{k-1}+\Delta m_k}.
\end{align}
The products $(\rho T)_f$ and $(\rho^2)_f$ are interpolated products,
not products of separately averaged quantities. The inward turbulent flux at
the physical inner boundary and the surface turbulent flux are zero.
The existing forced-nonturbulent masks are retained. The flux residual
differentiates these expressions directly. For example,
\begin{equation}
 \delta L_r=\mathrm{Lrad\_coeff}\,(\delta\nabla_{L,\rm eff}+\delta Y)
              +\nabla\,\delta\mathrm{Lrad\_coeff}.
\end{equation}
At fixed structure and nonnegative w, convective luminosity has a
nonnegative Y derivative and the radiative term has a positive Y derivative.
This local property does not prove convergence of the coupled stellar solve.
TDC retains its own convection relation and viscosity routines.

\subsection{Face source and damping}
The one equation source uses the face EOS and local PII:
\begin{equation}
 S_f=(w_f+w_{\rm seed})\frac{\mathrm{PII}_f}{Hp_f}
 \frac{P_f\chi_{T,f}}{\rho_f\chi_{\rho,f}C_{P,f}}.
 \label{eq:rsp2source}
\end{equation}
With $x_D=(8/3)\sqrt{2/3}$ and $x_R=2\sqrt3$,
\begin{align}
 D_f&=\alpha_dx_D\frac{w_f^3}{\Lambda_f},\label{eq:rsp2damp}\\
 D_{r,f}&=4\sigma_{\rm SB}(\alpha_rx_R)^2
 \frac{w_f^2T_f^3}{\rho_f^2C_{P,f}\kappa_f\Lambda_f^2}.
 \label{eq:rsp2raddamp}
\end{align}
These are \code{compute_Source}, \code{compute_D}, and \code{compute_Dr}.
LNA calls those same routines and uses $C=S-D-D_r$.
They use the public face length functions; the cell length is reserved
for cell-centered stress coefficients and turbulent transport.

\section{Eddy viscosity: both grids and both work forms}
\label{app:viscosity}

TDC viscosity is implemented in \code{tdc_hydro.f90}; RSP2 viscosity is in
\code{hydro_rsp2.f90}.
The routines remain separate. Both use face turbulent velocities with their
own closure normalizations and time weights; the public length functions
are shared.
The static LNA uses the equilibrium stress coefficients described below.

\subsection{Face velocity: cell stress}

For face velocity v, both models use $\Lambda_{\mathrm{cell},k}$, but the
stress velocity is $(w_k+w_{k+1})/2$ in RSP2 and
$(v_{c,k}+v_{c,k+1})/(2\sqrt{2/3})$ in TDC. Both use half the available
face value in the last cell with a zero inner moment boundary. The cell stress is
\begin{equation}
 Chi_k=\frac{16\pi}{3}\frac{\alpha_m\rho_k^2}{\Delta m_k}
       \frac{r_k^6+r_{k+1}^6}{2}\Lambda_{\mathrm{cell},k}w_{\mathrm{stress},k}
       \left(\frac{v_k}{r_k}-\frac{v_{k+1}}{r_{k+1}}\right).
\end{equation}
Here $w_{\mathrm{stress}}$ denotes the placement just specified; it is not a
new solver variable. Since the background strain is zero, only the velocity
variation contributes to $\delta Chi$ at first order. Then
\begin{equation}
 \delta Uq_k=\frac{4\pi}{r_k\Delta m_{\mathrm{face},k}}
                  (\delta Chi_{k-1}-\delta Chi_k).
\end{equation}
The dual mass is $(\Delta m_{k-1}+\Delta m_k)/2$ for $k>1$ and
$\Delta m_1/2$ at the surface. The exterior stress is zero. Nonlinear
RSP2 also retains the native mass-correction factors in this dual mass;
LNA does not support those factors.
\textit{Code:} \code{compute_Chi_cell}, \code{compute_Uq_face}, and the
LNA \code{static_eddy_Chi_cell_for_star_LNA}/\code{static_eddy_Uq_face_for_star_LNA}.

\subsection{Cell velocity: face stress}

For \code{u_flag}, RSP2 uses its local face w and TDC uses its stored
face $v_c/\sqrt{2/3}$. The face stress is
\begin{equation}
 Chi_{f,k}=\frac{16\pi}{3}
 \frac{\alpha_m\rho_{f,k}^2r_k^6\Lambda_{f,k}w_{f,k}}
      {\Delta m_{\mathrm{face},k}}
 \left(\frac{u_{k-1}}{r_{\mathrm{mid,start},k-1}}
       -\frac{u_k}{r_{\mathrm{mid,start},k}}\right).
\end{equation}
The nonlinear force is
\begin{equation}
 \Delta m_k Uq_k=\frac{4\pi}{r_{\mathrm{mid,start},k}}
                         (Chi_{f,k}-Chi_{f,k+1}).
\end{equation}
RSP2 \code{compute_Uq_dm_cell} returns this force. The Riemann momentum
row divides it by dm exactly once; it is not added to the contact velocity.
On the static background, the corresponding LNA term is
\begin{equation}
 \delta Uq_k=\frac{4\pi}{r_{\mathrm{mid},k}\Delta m_k}
                         (\delta Chi_{f,k}-\delta Chi_{f,k+1}).
 \label{eq:uqcell}
\end{equation}
\textit{Code:} \code{static_eddy_Chi_face_for_star_LNA} and
\code{Uq_cell_for_star_LNA}. Density, length, and w perturbations multiply
zero background strain, so their products have no first-order contribution.

The surface face stress is zero. A positive \code{R_center} permits an inner
boundary stress when \path{TDC_include_inner_boundary_eddy_viscosity} is true.
RSP2 also requires an empty inner forced-nonturbulent region.
It uses $R_{\rm center}$, $\rho_N$, $\Lambda_{f,N}$, and mass
$\Delta m_N/2$. Its turbulent velocity is the innermost stored face w for RSP2,
and $v_{c,N}/\sqrt{2/3}$ for TDC. The strain is
$u_N/r_{\mathrm{mid,start},N}-v_{\rm center}/R_{\rm center}$; the prescribed
inner velocity has zero perturbation. A temperature-selected LNA domain does
not create this physical boundary.

\subsection{Nonlinear heating and the work identity}

Use the same stress in momentum and heating. On the v grid,
\begin{equation}
 Eq_{\mathrm{stress},k}=\frac{4\pi Chi_k}{\Delta m_k}
 \left(\frac{v_k^{\rm work}}{r_k^{\rm work}}
       -\frac{v_{k+1}^{\rm work}}{r_{k+1}^{\rm work}}\right).
\end{equation}
For RSP2 this stress heating is allocated to the face energy volumes:
\begin{equation}
 Eq_{f,k}=w_k\frac{\Delta m_{k-1}(Eq/w)_{c,k-1}+\Delta m_k(Eq/w)_{c,k}}
 {\Delta m_{k-1}+\Delta m_k}.
\end{equation}
The coefficient $(Eq/w)_c$ is evaluated before multiplying by the cell
stress velocity. The gas row retains $Eq_{c,k}=Eq_{\mathrm{stress},k}$,
as in TDC. The cell and face heating integrals agree. No division by a
vanishing cell stress velocity or projection back into cells is required.
On the u grid, first evaluate the face-specific heating,
\begin{align}
 Eq_{f,k}&=\frac{4\pi Chi_{f,k}}{\Delta m_{\mathrm{face},k}}
 \left(\frac{u_{k-1}^{\rm work}}{r_{\mathrm{mid,start},k-1}}
       -\frac{u_k^{\rm work}}{r_{\mathrm{mid,start},k}}\right),\\
 Eq_k&=\tfrac12(Eq_{f,k}+Eq_{f,k+1}).
\end{align}
The half-face allocation sums to the full dual mass on an unequal-mass mesh.
The inner-boundary face has its own half-cell mass and prescribed velocity.
RSP2 deposits its inner-wall heating in the last evolved face volume,
keeping the prescribed inner moment zero.
Summation by parts gives
\begin{align}
 \sum_k\Delta m_{\mathrm{face},k}v_k^{\rm work}Uq_k
  +\sum_k\Delta m_k Eq_k
  &=-4\pi Chi_N v_{\rm center}/R_{\rm center},&&v\text{ grid},\\
 \sum_k\Delta m_k u_k^{\rm work}Uq_k
  +\sum_k\Delta m_k Eq_k
  &=-4\pi Chi_{f,N+1}v_{\rm center}/R_{\rm center},&&u\text{ grid}.
\end{align}
The right sides vanish for zero boundary stress or velocity. For a full center,
these expressions are understood by the zero-stress limit, not division by zero.
RSP2 masks stresses at the selected forced regions and retains their force
divergence at the interface. TDC retains its existing additional momentum
masks; the unrestricted identity is not asserted across those masks.

The four \code{dedt} cases have distinct energy-source terms:
\begin{center}
\small
\begin{tabular}{@{}llll@{}}
\toprule
Grid & Local PdV & Energy source from viscosity & Explicit mechanical inertia\\
\midrule
v & true  & $Eq_k$ & none\\
v & false & $Eq_k+\tfrac12(\bar v_kUq_k+\bar v_{k+1}Uq_{k+1})$ & $K+\Phi$\\
u & true  & $Eq_k$ & none\\
u & false & $Eq_k+\bar u_kUq_k$ & $K+\Phi$\\
\bottomrule
\end{tabular}
\end{center}
Here $\bar v=(v+v_{\rm start})/2$ and similarly for u. The v-grid power
also includes the native cell mass-correction factor when enabled in nonlinear
hydro. Conservative work always uses these averaged velocities, including
when optional velocity time centering is off. For local PdV work, heating uses
the average only if velocity time centering is on; otherwise it uses the current
velocity. The latter is a current-velocity power identity, not exact
finite-step kinetic-energy change. RSP2 \code{eps_grav} retains Eq but no
explicit Uq power or mechanical inertia, for either setting of the work control.

TDC's explicit-momentum option uses old convective velocity in momentum and
current velocity in heating. This intentionally distinct choice is unchanged;
same-stress cancellation then has a lag difference. RSP2 uses current face w
in both terms and does not inherit that TDC option.

The face-w contribution to conservative inner-face work can reach zone
$k+2$ through viscosity or turbulent pressure. With \code{u_flag}, Riemann
contact velocity and pressure also carry the inner cell turbulent pressure
into the local work and momentum rows. The nonlinear Jacobian and LNA
retain these pressure derivatives before the face AD origin is shifted.
The ordinary three-zone solve remains sufficient when the extra entries
vanish, including the \code{v_flag} local-work case. Viscous heating itself
uses the direct cell expression above on that grid.

\subsection{Static energy order}

On a zero-velocity background, $Uq=O(\delta)$ and
$Eq=O(\delta^2)$. Thus
\begin{equation}
 \delta(vUq)=\delta(uUq)=0,\qquad \delta Eq=0.
\end{equation}
\code{turbulent_viscous_heating_for_star_LNA} returns zero, while the
first-order Uq response remains in momentum. The nonlinear work comparisons
above therefore explain conservation without introducing additional LNA rows.

\subsection{Optional three equation closure and discretization}
\label{app:rsp2-moments}

\paragraph{Pi and PII.}
PII has units of specific entropy. Pi has units of specific entropy times
velocity. The old luminosity is
$L_{c,f}=\mathcal A_f(\rho T)_fw_f\mathrm{PII}_f$, whereas the new closure is
\begin{align}
 L_{c,f}&=\mathcal A_f\rho_fT_f\mathrm{Pi}_f,\\
 \delta L_{c,f}&=\mathcal A_f\rho_fT_f\delta\mathrm{Pi}_f
 +\mathrm{Pi}_f\delta(\mathcal A_f\rho_fT_f).
\end{align}
Pi replaces the algebraic heat transport closure. PII is not another solver
variable in this mode. Its retained diagnostic is $\mathrm{Pi}_f/w_f$ when
$w_f>0$, and zero when both vanish. When creating missing moments at an
unforced, unstable face with positive old $L_c$, the initial Pi is that
luminosity divided by the new $\mathcal A_f\rho_fT_f$; different face
interpolations mean this need not equal the old $w_f\mathrm{PII}_f$.
The initial variance is $3\mathrm{Pi}_f^2/(2e_{t,f})$ for positive face energy.
Other faces start with zero Pi and Phi, retaining their turbulent energy.
This initial-state assumption preserves the imported heat flux only at the
selected unstable faces; it does not assert moment equilibrium or modify
an existing three equation restart state.

\paragraph{Face state, source and cooling.}
The new closure uses the common reconstructed face EOS when enabled, or
the selected RSP2 face weights for each EOS quantity. Buoyancy uses the
resolved cell pressure difference:
\begin{equation}
 \left(-\frac{\partial_r P}{\rho}\right)_f
 =\frac{\mathcal A_f(P_k-P_{k-1})}{(\Delta m_{k-1}+\Delta m_k)/2}.
\end{equation}
The face energy source is local to its independent covariance:
\begin{equation}
 S_f=\left(-\frac{\partial_r P}{\rho}\right)_f
 \frac{\chi_{T,f}}{\chi_{\rho,f}C_{P,f}}\mathrm{Pi}_f.
\end{equation}
Its derivative retains all buoyancy and Pi terms. No averaging of Pi into
cells, reconstruction back to a face, or covariance cap enters the source.
The cooling rate and two new coefficients are
\begin{equation}
 \tau_{{\rm rad},f}^{-1}=
 \frac{4\sigma_{\rm SB}(\alpha_Rx_R)^2T_f^3}
 {C_{P,f}\kappa_f\rho_f^2\Lambda_f^2},\qquad
 \alpha_{\mathrm{Pi}}=6\sqrt{2/3},\quad
 \alpha_{\mathrm{Phi}}=4\sqrt{2/3}.
\end{equation}
The alpha symbols are the effective coefficients. In source they are
\code{RSP2_alfa_pi*x_ALFAPI} and \code{RSP2_alfa_phi*x_ALFAPHI};
the two controls default to one, and the fixed factors above are stored
in \code{hydro_rsp2}. These are the local MLT calibration listed in
\href{https://arxiv.org/html/2604.06151v1#S2}{Braun et al. (2026), section 2},
not that paper's fitted solar coefficients. The factor of two
in variance cooling follows from evolving the full $\langle s'^2\rangle$.

The normalization of the two existing controls is
\begin{equation}
 \gamma_r=2\sqrt3\,\code{RSP2_alfar},\qquad
 \alpha_\omega=\code{RSP2_alfat}.
\end{equation}
Braun et al. (2026), Table 1, list $\gamma_r=3.46$ and
$\alpha_\omega=0.25$. These correspond to \code{RSP2_alfar=1d0}
and \code{RSP2_alfat=0.25d0}; the present defaults are both zero.
The cooling rate scales as \code{RSP2_alfar} squared. Its physical role changes
with the model: the default model loses face turbulent energy through
$D_r=e_t/\tau_{\rm rad}$, while the three equation model cools face Pi and Phi
at rates $1/\tau_{\rm rad}$ and $2/\tau_{\rm rad}$ and has no direct $D_r$.
Alfat controls kinetic energy diffusion with coefficient
$\alpha_\omega\Lambda\sqrt{e_t}$ in both models. There is no extra
$\sqrt{2/3}$ in its coefficient and no added transport of Pi or Phi.
Matching coefficients does not make the face interpolation, mixing length
prescription or complete models identical.

\paragraph{Entropy driving and the temperature row.}
\code{get_rsp2_thermal_gradient} maps the selected temperature relation to
its radial thermal differential. At uniform composition,
\begin{equation}
 (-\partial_r s)_f=C_{P,f}\left[
    (-\partial_r\ln T)_{{\rm row},f}
       -\nabla_{{\rm ad},f}(-\partial_r\ln P)_{{\rm actual},f}\right].
\end{equation}
The QHSE and radiation pressure forms use the EOS differential
\begin{equation}
 (-\partial_r\ln P)_{{\rm actual},f}
 =\frac{\mathcal A_f\rho_f(P_k-P_{k-1})}{\Delta m_{c,f}P_f}.
\end{equation}
The logarithmic row uses its actual logarithmic pressure difference instead.
For QHSE, the coefficient multiplying $\nabla_T$ is
$-\mathcal A_f\rho_f h_fT_{*,f}/T_f$; for radiation pressure it is
$\rho_f\kappa_fL_{r,\mathrm{coeff},f}/
(4c\mathcal A_f\lambda_fP_{{\rm rad},f})$.
The logarithmic row uses
$\mathcal A_f\rho_f\ln(P_k/P_{k-1})/\Delta m_{c,f}$.
The radiation coefficient, opacity floor and boundary spacing are shared
with the actual temperature row.

For the first two forms, define the neutral gradient using the selected
reference pressure differential and the temperature coefficient just given.
Dynamic gradL uses the actual pressure differential; static gradL uses the
HSE pressure differential. The algebraically equivalent evaluation is
\begin{equation}
\begin{aligned}
 (-\partial_r s)_f=C_{P,f}\big[&
 (-\partial_r\ln T)_{\nabla_T=1,f}\,Y_f
 +B_f(-\partial_r\ln P)_{{\rm ref},f}\\
 &+\nabla_{{\rm ad},f}\{
 (-\partial_r\ln P)_{{\rm ref},f}
 -(-\partial_r\ln P)_{{\rm actual},f}\}\big].
\end{aligned}
\end{equation}
Here $B_f$ is the existing composition contribution to gradL. The dynamic
branch cancels the last difference analytically. The logarithmic form uses
$\nabla_L=\nabla_{\rm ad}+B$ and
$(-\partial_r s)=C_P(-\partial_r\ln P)_{\rm log}(Y+B)$.
All current structure derivatives are retained; the existing composition
coefficient is held fixed in thermal AD. These algebraic cancellations
change neither implicitness nor luminosity time weights. Nonlinear
\code{constant_L} retains a measured entropy differential; LNA still excludes
that temperature row.

\paragraph{Mean velocity strain.}
For cell velocity, the face strain is
$\mathcal A_f\rho_f(u_{k-1}-u_k)/\Delta m_{c,f}$.
For face velocity it is
\begin{equation}
 (\partial_r v)_f=a_f\frac{v_k-v_{k+1}}{r_k-r_{k+1}}
 +(1-a_f)\frac{v_{k-1}-v_k}{r_{k-1}-r_k}.
\end{equation}
The latter is exact for an affine radial velocity on an unequal grid.
It is a deformation term in Pi evolution, with no new viscosity coefficient.

\paragraph{Nonlinear residuals and normalization.}
Each moment row is its current value minus its accepted value minus
$\Delta t$ times its current rhs, divided by a fixed start scale.
Using the same positive luminosity scale $L_{\rm scale}$ as the flux row,
\begin{align}
 \mathrm{Pi}_{\rm ref}&=L_{\rm scale}/(\mathcal A\rho T),
 &v_{\rm ref}&=[L_{\rm scale}/(\mathcal A\rho)]^{1/3},\\
 \mathrm{Pi}_{\rm scale}&=\max(|\mathrm{Pi}_{\rm start}|,\mathrm{Pi}_{\rm ref}),
 &\mathrm{Phi}_{\rm scale}&=\max(\mathrm{Phi}_{\rm start},
                \tfrac32[\mathrm{Pi}_{\rm ref}/v_{\rm ref}]^2).
\end{align}
These also scale the new solver columns. They impose no bound on Pi.
The one-equation RSP2 model retains the original normalized turbulent energy
residual for nonzero accepted energy. The three-equation model multiplies
that residual by $\max(1,c_s/w)$. For $0<w<c_s$, known factors of $w$
are canceled term by term before differentiation. In particular,
$(w^2-w_{\rm start}^2)/w=w-w_{\rm start}^2/w$ retains its derivative
$1+(w_{\rm start}/w)^2$, and the explicit numerator $w$ cancels from
the buoyancy source. Face energy and all remaining state derivatives are
retained. Face transport and $u$ grid stress use the exact ratio
$[a w+(1-a)w_{\rm neighbor}]/w=a+(1-a)w_{\rm neighbor}/w$;
accepted energy, pressure and flux terms remain divided by current $w$.
Remaining quotients use $q=x/w$, $\delta q=(\delta x-q\delta w)/w$,
avoiding an intermediate $1/w^2$ that can overflow at tiny $w$.
Physical equations and time weights are unchanged. At a stationary positive
energy background, scaling both stiffness and inertia preserves the LNA
eigenvalues. The zero energy branch is selected as described with the
moment equations.
The companion RSP3 manuscript gives the complete divided row.
At a factorizable zero accepted energy, the analytic canceled row selects the zero
or positive branch in either model. The three-equation divided row retains
a positive trial value when accepted energy or surviving moments require it.
The energy test includes exact underflow of $w_{\rm start}^2$ to zero;
positive representable history is retained. No denominator floor is used.
The positive multiplier is at least one, so meeting the scaled residual
tolerance also meets the original normalized energy tolerance.

The local isotropic closure does not guarantee
$\mathrm{Pi}_f^2\leq(2/3)e_{t,f}\mathrm{Phi}_f$.
Independent interpolation and kinetic energy transport do not establish that
bound either. The covariance excess profile reports the difference; it does
not limit heat flux. The two new moment perturbations are included in LNA
eigenfunction output, and setup reports their dimensional equilibrium defects.

\section{Mode and work conventions}

Finite positive frequency roots are sorted by increasing $\omega$.  The
selected list applies the frequency floor and the existing growth controls
to
\begin{equation}
  g_{\ln K}=\code{logKE_per_cycle}
  =\frac{4\pi\gamma}{\omega}.
  \label{eq:logke}
\end{equation}
This is logarithmic kinetic energy growth per cycle, not a generic GYRE
$\eta$.  Derived output quantities are
\begin{align}
  \frac{\Delta K}{K} &=\exp(g_{\ln K})-1,\\
  \mathrm{GREKM} &=2\tanh(g_{\ln K}/2),\\
  \frac{\Delta a}{a} &=\exp(g_{\ln K}/2)-1.
\end{align}
The old \code{star_LNA_*_eta} identifiers are retained for inlist
compatibility, but comments and files state that they filter
\code{logKE_per_cycle}.

The optional velocity initialization combines the selected complex velocity
eigenfunctions.  For mode $j$, the surface-phase normalization is
\begin{equation}
  \widehat{\delta v}_{j,k}=\frac{\delta v_{j,k}}{\delta v_{j,1}}.
\end{equation}
The requested fractions are applied to these normalized eigenfunctions, and
the real velocity is scaled to \code{star_LNA_kick_vsurf_km_per_sec} at the
surface.  This preserves the computed velocity phase and amplitude instead of
constructing a velocity profile from displacement amplitudes.
With \code{u_flag}, the active eigenfunction is cell velocity. The kick
updates its solver state and rebuilds the Riemann face velocities and their
start values. Only the velocity state is initialized.

Pressure work is reported with the RSP phase convention
\begin{equation}
  W_{P,k}=-\pi\Delta m_k
  \operatorname{Im}\left(\delta P_k^*\delta\nu_k\right),
\end{equation}
where $\nu=1/\rho$ is specific volume and
$\delta\nu_k=-\delta\ln\rho_k/\rho_k$. The result is divided by the modal
kinetic energy normalization. EOS and RSP2 turbulent
pressure perturbations use the same closure as the operator.  The eddy viscosity
and luminosity terms have not been compared with RSP/RSP2 work output.  The
reported \code{total_work} is an arithmetic sum of six diagnostics: EOS
pressure, turbulent pressure, eddy viscosity, and radiative, convective,
and turbulent luminosity work. These six quantities have not been derived
and validated as one discrete work balance against Equation~\eqref{eq:logke}.
Pressure work and luminosity-based thermal diagnostics need not represent
independent energy exchanges, so their simple addition is not an established
decomposition of net modal driving. A consistent work identity should provide
an independent check of the eigenvalue growth. This limitation of the
postprocessing does not itself establish an error in the eigenvalue.

Routine \code{eddy_viscous_work_for_star_LNA} uses the acceleration from the
same momentum operator. Before division by modal kinetic energy,
\begin{equation}
  W_{q,k}=\frac{\pi\Delta m_k}{|\omega|}
  \begin{cases}
    \operatorname{Re}(\delta u_k^*\delta U_{q,k}),&\code{u_flag},\\
    \tfrac12\operatorname{Re}(\delta v_k^*\delta U_{q,k}
       +\delta v_{k+1}^*\delta U_{q,k+1}),&\text{face velocity}.
  \end{cases}
\end{equation}
Perturbations below the active domain are zero. The normalization is
$K=\tfrac12\sum_k\Delta m_{v,k}|\delta v_{{\rm rad},k}|^2$, using cell
masses for $u$ and \code{star_LNA_dm_face} for $v$.


\section{Schema 130 TDC background export}

When \code{gyre_data_schema = 130} and
\code{gyre_write_tdc_lna_background = .true.},
\sourcefile{star/private/pulse_gyre.f90} writes the standard schema-120 columns
1 through 19 followed by the columns in Table~\ref{tab:schema130}.  The exporter
writes background data only; it does not run GYRE or define GYRE's perturbation
equations.

Schema 130 does not carry the harmonic TDC mixing length. When
\path{harmonic_dissipation_length_beta} is positive, MESA writes a warning and
still exports the legacy \code{mixing_length_alpha} and \code{Hp_face} fields.
The legacy length represented by these fields is
\begin{equation}
  \Lambda_{\mathrm{GYRE}}=\alpha_{\mathrm{MLT}}H_P,
\end{equation}
which differs from the harmonic $\Lambda$ used by MESA. This describes the
MESA exporter; changes in a separate GYRE checkout do not change the STAR LNA
operator or this file schema.

\begin{table}[htbp]
\centering
\footnotesize
\caption{Additional schema 130 point columns.}\label{tab:schema130}
\begin{tabular}{@{}rll@{}}
\toprule
Column & Source quantity & Meaning \\
\midrule
20 & \code{L_conv} & background convective luminosity $L_{c,0}$ \\
21 & \code{mlt_vc/sqrt_2_div_3} & TDC amplitude $A_0$ \\
22 & \code{i_tdc_d_conv_h} & inferred diffusivity $D_{\mathrm{conv}}$ \\
23 & \code{Hp_face} & face pressure scale height \\
24 & \code{mixing_length_alpha} & $\alpha_{\mathrm{MLT}}$ \\
25 & \code{Cp_face} & face specific heat \\
26 & \code{chiT_face} & $\chi_T$ \\
27 & \code{chiRho_face} & $\chi_\rho$ \\
28 & \code{gradL} & Ledoux gradient \\
29 & \code{gradT} & actual temperature gradient \\
30 through 35 & TDC controls & $\alpha_C,\alpha_S,\alpha_D,\alpha_R,\alpha_{Pt},\alpha_M$ \\
36 & logical as real & include the MLT correction \\
37 & \code{mlt_Pturb_factor} & MLT turbulent pressure factor \\
38 & inferred & background MLT turbulent pressure \\
\bottomrule
\end{tabular}
\end{table}

The exported diffusivity is inferred by \code{store_tdc_lna_point_env} from
\begin{equation}
  F_{\mathrm{conv}}
  =\rho TC_P D_{\mathrm{conv}}
  \frac{\nabla_T-\nabla_L}{H_P},
\end{equation}
so
\begin{equation}
  D_{\mathrm{conv}}
  =\frac{L_{\mathrm{conv}}H_P}
  {4\pi r^2\rho TC_P(\nabla_T-\nabla_L)}
\end{equation}
when the denominator and $L_{\mathrm{conv}}$ are positive. Its units are
$\mathrm{cm^2\,s^{-1}}$; the earlier label $D_{\mathrm{conv}}H_P$ was
inconsistent with this equation and the exporter. The exported
turbulent pressure is
\begin{equation}
  P_{\mathrm{t,MLT},0}
  =f_{\mathrm{Pt}}\frac{\rho_fv_c^2}{3}.
\end{equation}


\section{Audit files and restrictions}

With \code{star_LNA_write_matrix_summary = .true.}, the source writes
\begin{itemize}
\item \code{<prefix>_matrix_summary.data}, including matrix dimensions, nonzero
      counts, norm ranges, residuals, and equation counts;
\item \code{<prefix>_row_structure.data}, including the equation name, nonzero
      counts, and dominant A/B variable for every row;
\item \code{<prefix>_rsp2_term_audit.data} on RSP2 backgrounds, including
      PII, $L_c$, $L_t$, source, damping, radiative damping, turbulent pressure
      work coefficient, divergences, RHS, and inertia; and
\item \code{<prefix>_tdc_face_audit.data} for active TDC, including active and
      stored face states, background convection quantities, closure
      luminosities, velocity RHS, and inertia.
\end{itemize}
When the full-pencil residual rejects a root that otherwise passes the mode
heuristics, the terminal output identifies the first such period and its
maximum-residual row, zone, and row variable.

The source has the following restrictions:\par
\begin{itemize}
\setlength{\itemsep}{1pt}
\setlength{\parsep}{0pt}
\item no rotation, RTI, mass correction, or velocity drag;
\item no \path{use_compression_outer_BC}, \code{constant_L}, user
      \code{other_*};
\item \code{eps_grav} is supported only for RSP2; no nonzero \code{mstar_dot};
\item active RSP2 or MLT turbulent pressure with \code{u_flag} requires
      \path{star_LNA_perturb_turbulent_pressure} to be true;
\item the optional initial velocity kick requires \code{v_flag} or
      \code{u_flag}; the \code{u_flag} path rebuilds the Riemann face velocity
      and initializes its start-of-step state;
\item RSP2 requires \code{star_LNA_convection_treatment = 'perturbed'};
\item no derivative of the nonlinear enthalpy flux limiter in the TDC
      operator;
\item dense all roots scaling with model zone count;
\item the six-column \code{total_work} diagnostic is not a derived discrete
      work balance against the eigenvalue growth.
\end{itemize}

\code{check_star_LNA_model} rejects incompatible controls. Setup output lists
accepted controls that alter the linearization. The existing three-zone AD
stencil can drop farther-neighbor derivatives after shifts; the assembled
operator retains that supported stencil, rather than an unrestricted
full-band derivative of every composite closure.


\section{Surface boundary details and source map}
\label{app:surface-details}

\subsection{Atmospheric boundary offsets}

The ordinary atmospheric temperature condition uses
\begin{align}
  T_{\rm bc}&=T_{\rm atm}+\Delta T_0,\qquad
  \Delta T_0=\Delta P_0\frac{\nabla_1T_1}{P_{{\rm eos},1}},\\
  \Delta P_0&=\frac{G_1m_{g,1}\Delta m_1}{8\pi r_1^4}.
\end{align}
The offset is retained for both choices of \code{use_momentum_outer_BC}
and differentiated with AD. At fixed mass and composition,
\begin{align}
  \delta\Delta T_0
  &=\frac{\Delta P_0T_1}{P_{{\rm eos},1}}\,\delta\nabla_1
    +\Delta T_0\left[-4\delta\ln r_1
       +(1-\chi_{T,1})\delta\ln T_1-\chi_{\rho,1}\delta\ln\rho_1\right],\\
  \delta\ln T_{\rm bc}
  &=\frac{T_{\rm atm}\delta\ln T_{\rm atm}+\delta\Delta T_0}
          {T_{\rm atm}+\Delta T_0}.
\end{align}
The atmosphere supplies $\delta\ln T_{\rm atm}$, including its supported
opacity, luminosity, and radius responses. This is the boundary variation
used in the surface-temperature row in Section~\ref{sec:thermal-boundary}.

When \code{use_momentum_outer_BC} and
\path{floor_momentum_outer_BC_at_Prad} are true, the momentum boundary uses
\begin{equation}
  P_{\rm bc}=\max\left(P_{\rm bc}^{\rm raw},\,aT_{\rm atm}^4/3\right).
\end{equation}
On the floored branch, $\delta\ln P_{\rm bc}=4\delta\ln T_{\rm atm}$;
otherwise the pre-floor pressure response is retained. Momentum pressure
remains at the outer face, while the temperature offset places the
temperature at the center of cell 1. The luminosity boundary in
Section~\ref{sec:thermal-boundary} uses no atmospheric temperature offset.


\paragraph{Nonlinear surface-row details.}
The RSP2-specific switch defaults to true; the shared switch defaults to false.
With both off, the normal atmosphere temperature condition applies. The
separate \code{constant_L} override is rejected by STAR LNA.
The nonlinear RSP2 row, \code{do1_rsp2_L_eqn}, uses current area
$4\pi r_1^2$ and normalizes by
$\max(|L_{r,1}|,|L_1|,1\,\mathrm{erg\,s^{-1}})$; $L_{c,1}=L_{t,1}=0$.
The shared \code{set_RSP_Lsurf_BC} uses \code{maxval(L_start)} and, when
velocity time centering is active, the area
\begin{equation}
 \mathcal A_1^{\rm hydro}=\frac{4\pi}{3}
 \left(r_1^2+r_1r_{1,\rm start}+r_{1,\rm start}^2\right).
\end{equation}
Thus the two nonlinear rows can differ during a step. Their static boundary
values agree; the continuous LNA uses current area and the linearized row
above, without finite-step centering. The additional RSP2 surface row remains
$Y_1=0$ and does not duplicate this luminosity boundary.

\subsection{Source files}

The row routines cited in the main text are in the following files:

\begin{itemize}
\setlength{\itemsep}{1pt}
\setlength{\parsep}{0pt}
\item \sourcefile{star/private/star_LNA.f90}, for the public solve flow and row
      assembly order;
\item \sourcefile{star/private/star_LNA_support.f90}, for rows, matrix
      construction, algebraic elimination, mode output, work output, and
      audits;
\item \sourcefile{star/private/star_LNA_turbulence_closures.f90}, for RSP2 and
      TDC pressure, luminosity, turbulent energy, and eddy viscosity closures;
\item \sourcefile{star/private/hydro_rsp2.f90}, for the independent-Y residual,
      PII closure, shared-length calls, and RSP2 viscosity;
\item \sourcefile{star/private/tdc_hydro.f90}, for the shared lengths and
      unchanged TDC viscosity;
\item \sourcefile{star/private/turb_support.f90}, for the shared dynamical
      neutral gradient;
\item \sourcefile{star/private/hydro_temperature.f90} and
      \sourcefile{star/private/hydro_eqns.f90}, for nonlinear temperature rows
      and surface-boundary selection;
\item \sourcefile{star/private/hydro_energy.f90} and
      \sourcefile{star/private/eps_grav.f90}, for nonlinear energy forms;
\item \sourcefile{star/private/hydro_riemann.f90} and
      \sourcefile{star/private/auto_diff_support.f90}, for cell momentum
      and reconstructed face variables;
\item \sourcefile{star/private/star_utils.f90}, for discrete mechanical
      energies;
\item \sourcefile{turb/public/turb.f90}, routine \code{set_TDC_LNA}, for the
      local TDC relation; and
\item \sourcefile{star/private/pulse_gyre.f90}, for the schema 130 TDC
      background export.
\end{itemize}


\end{document}
